\documentclass[11pt]{article}

    \usepackage[breakable]{tcolorbox}
    \usepackage{parskip} % Stop auto-indenting (to mimic markdown behaviour)
    

    % Basic figure setup, for now with no caption control since it's done
    % automatically by Pandoc (which extracts ![](path) syntax from Markdown).
    \usepackage{graphicx}
    % Maintain compatibility with old templates. Remove in nbconvert 6.0
    \let\Oldincludegraphics\includegraphics
    % Ensure that by default, figures have no caption (until we provide a
    % proper Figure object with a Caption API and a way to capture that
    % in the conversion process - todo).
    \usepackage{caption}
    \DeclareCaptionFormat{nocaption}{}
    \captionsetup{format=nocaption,aboveskip=0pt,belowskip=0pt}

    \usepackage{float}
    \floatplacement{figure}{H} % forces figures to be placed at the correct location
    \usepackage{xcolor} % Allow colors to be defined
    \usepackage{enumerate} % Needed for markdown enumerations to work
    \usepackage{geometry} % Used to adjust the document margins
    \usepackage{amsmath} % Equations
    \usepackage{amssymb} % Equations
    \usepackage{textcomp} % defines textquotesingle
    % Hack from http://tex.stackexchange.com/a/47451/13684:
    \AtBeginDocument{%
        \def\PYZsq{\textquotesingle}% Upright quotes in Pygmentized code
    }
    \usepackage{upquote} % Upright quotes for verbatim code
    \usepackage{eurosym} % defines \euro

    \usepackage{iftex}
    \ifPDFTeX
        \usepackage[T1]{fontenc}
        \IfFileExists{alphabeta.sty}{
              \usepackage{alphabeta}
          }{
              \usepackage[mathletters]{ucs}
              \usepackage[utf8x]{inputenc}
          }
    \else
        \usepackage{fontspec}
        \usepackage{unicode-math}
    \fi

    \usepackage{fancyvrb} % verbatim replacement that allows latex
    \usepackage[Export]{adjustbox} % Used to constrain images to a maximum size
    \adjustboxset{max size={0.9\linewidth}{0.9\paperheight}}

    % The hyperref package gives us a pdf with properly built
    % internal navigation ('pdf bookmarks' for the table of contents,
    % internal cross-reference links, web links for URLs, etc.)
    \usepackage{hyperref}
    % The default LaTeX title has an obnoxious amount of whitespace. By default,
    % titling removes some of it. It also provides customization options.
    \usepackage{titling}
    \usepackage{longtable} % longtable support required by pandoc >1.10
    \usepackage{booktabs}  % table support for pandoc > 1.12.2
    \usepackage{array}     % table support for pandoc >= 2.11.3
    \usepackage{calc}      % table minipage width calculation for pandoc >= 2.11.1
    \usepackage[inline]{enumitem} % IRkernel/repr support (it uses the enumerate* environment)
    \usepackage[normalem]{ulem} % ulem is needed to support strikethroughs (\sout)
                                % normalem makes italics be italics, not underlines
    \usepackage{mathrsfs}
    

    
    % Colors for the hyperref package
    \definecolor{urlcolor}{rgb}{0,.145,.698}
    \definecolor{linkcolor}{rgb}{.71,0.21,0.01}
    \definecolor{citecolor}{rgb}{.12,.54,.11}

    % ANSI colors
    \definecolor{ansi-black}{HTML}{3E424D}
    \definecolor{ansi-black-intense}{HTML}{282C36}
    \definecolor{ansi-red}{HTML}{E75C58}
    \definecolor{ansi-red-intense}{HTML}{B22B31}
    \definecolor{ansi-green}{HTML}{00A250}
    \definecolor{ansi-green-intense}{HTML}{007427}
    \definecolor{ansi-yellow}{HTML}{DDB62B}
    \definecolor{ansi-yellow-intense}{HTML}{B27D12}
    \definecolor{ansi-blue}{HTML}{208FFB}
    \definecolor{ansi-blue-intense}{HTML}{0065CA}
    \definecolor{ansi-magenta}{HTML}{D160C4}
    \definecolor{ansi-magenta-intense}{HTML}{A03196}
    \definecolor{ansi-cyan}{HTML}{60C6C8}
    \definecolor{ansi-cyan-intense}{HTML}{258F8F}
    \definecolor{ansi-white}{HTML}{C5C1B4}
    \definecolor{ansi-white-intense}{HTML}{A1A6B2}
    \definecolor{ansi-default-inverse-fg}{HTML}{FFFFFF}
    \definecolor{ansi-default-inverse-bg}{HTML}{000000}

    % common color for the border for error outputs.
    \definecolor{outerrorbackground}{HTML}{FFDFDF}

    % commands and environments needed by pandoc snippets
    % extracted from the output of `pandoc -s`
    \providecommand{\tightlist}{%
      \setlength{\itemsep}{0pt}\setlength{\parskip}{0pt}}
    \DefineVerbatimEnvironment{Highlighting}{Verbatim}{commandchars=\\\{\}}
    % Add ',fontsize=\small' for more characters per line
    \newenvironment{Shaded}{}{}
    \newcommand{\KeywordTok}[1]{\textcolor[rgb]{0.00,0.44,0.13}{\textbf{{#1}}}}
    \newcommand{\DataTypeTok}[1]{\textcolor[rgb]{0.56,0.13,0.00}{{#1}}}
    \newcommand{\DecValTok}[1]{\textcolor[rgb]{0.25,0.63,0.44}{{#1}}}
    \newcommand{\BaseNTok}[1]{\textcolor[rgb]{0.25,0.63,0.44}{{#1}}}
    \newcommand{\FloatTok}[1]{\textcolor[rgb]{0.25,0.63,0.44}{{#1}}}
    \newcommand{\CharTok}[1]{\textcolor[rgb]{0.25,0.44,0.63}{{#1}}}
    \newcommand{\StringTok}[1]{\textcolor[rgb]{0.25,0.44,0.63}{{#1}}}
    \newcommand{\CommentTok}[1]{\textcolor[rgb]{0.38,0.63,0.69}{\textit{{#1}}}}
    \newcommand{\OtherTok}[1]{\textcolor[rgb]{0.00,0.44,0.13}{{#1}}}
    \newcommand{\AlertTok}[1]{\textcolor[rgb]{1.00,0.00,0.00}{\textbf{{#1}}}}
    \newcommand{\FunctionTok}[1]{\textcolor[rgb]{0.02,0.16,0.49}{{#1}}}
    \newcommand{\RegionMarkerTok}[1]{{#1}}
    \newcommand{\ErrorTok}[1]{\textcolor[rgb]{1.00,0.00,0.00}{\textbf{{#1}}}}
    \newcommand{\NormalTok}[1]{{#1}}

    % Additional commands for more recent versions of Pandoc
    \newcommand{\ConstantTok}[1]{\textcolor[rgb]{0.53,0.00,0.00}{{#1}}}
    \newcommand{\SpecialCharTok}[1]{\textcolor[rgb]{0.25,0.44,0.63}{{#1}}}
    \newcommand{\VerbatimStringTok}[1]{\textcolor[rgb]{0.25,0.44,0.63}{{#1}}}
    \newcommand{\SpecialStringTok}[1]{\textcolor[rgb]{0.73,0.40,0.53}{{#1}}}
    \newcommand{\ImportTok}[1]{{#1}}
    \newcommand{\DocumentationTok}[1]{\textcolor[rgb]{0.73,0.13,0.13}{\textit{{#1}}}}
    \newcommand{\AnnotationTok}[1]{\textcolor[rgb]{0.38,0.63,0.69}{\textbf{\textit{{#1}}}}}
    \newcommand{\CommentVarTok}[1]{\textcolor[rgb]{0.38,0.63,0.69}{\textbf{\textit{{#1}}}}}
    \newcommand{\VariableTok}[1]{\textcolor[rgb]{0.10,0.09,0.49}{{#1}}}
    \newcommand{\ControlFlowTok}[1]{\textcolor[rgb]{0.00,0.44,0.13}{\textbf{{#1}}}}
    \newcommand{\OperatorTok}[1]{\textcolor[rgb]{0.40,0.40,0.40}{{#1}}}
    \newcommand{\BuiltInTok}[1]{{#1}}
    \newcommand{\ExtensionTok}[1]{{#1}}
    \newcommand{\PreprocessorTok}[1]{\textcolor[rgb]{0.74,0.48,0.00}{{#1}}}
    \newcommand{\AttributeTok}[1]{\textcolor[rgb]{0.49,0.56,0.16}{{#1}}}
    \newcommand{\InformationTok}[1]{\textcolor[rgb]{0.38,0.63,0.69}{\textbf{\textit{{#1}}}}}
    \newcommand{\WarningTok}[1]{\textcolor[rgb]{0.38,0.63,0.69}{\textbf{\textit{{#1}}}}}


    % Define a nice break command that doesn't care if a line doesn't already
    % exist.
    \def\br{\hspace*{\fill} \\* }
    % Math Jax compatibility definitions
    \def\gt{>}
    \def\lt{<}
    \let\Oldtex\TeX
    \let\Oldlatex\LaTeX
    \renewcommand{\TeX}{\textrm{\Oldtex}}
    \renewcommand{\LaTeX}{\textrm{\Oldlatex}}
    % Document parameters
    % Document title
    \title{Spectroscopy2}
    
    
    
    
    
% Pygments definitions
\makeatletter
\def\PY@reset{\let\PY@it=\relax \let\PY@bf=\relax%
    \let\PY@ul=\relax \let\PY@tc=\relax%
    \let\PY@bc=\relax \let\PY@ff=\relax}
\def\PY@tok#1{\csname PY@tok@#1\endcsname}
\def\PY@toks#1+{\ifx\relax#1\empty\else%
    \PY@tok{#1}\expandafter\PY@toks\fi}
\def\PY@do#1{\PY@bc{\PY@tc{\PY@ul{%
    \PY@it{\PY@bf{\PY@ff{#1}}}}}}}
\def\PY#1#2{\PY@reset\PY@toks#1+\relax+\PY@do{#2}}

\@namedef{PY@tok@w}{\def\PY@tc##1{\textcolor[rgb]{0.73,0.73,0.73}{##1}}}
\@namedef{PY@tok@c}{\let\PY@it=\textit\def\PY@tc##1{\textcolor[rgb]{0.24,0.48,0.48}{##1}}}
\@namedef{PY@tok@cp}{\def\PY@tc##1{\textcolor[rgb]{0.61,0.40,0.00}{##1}}}
\@namedef{PY@tok@k}{\let\PY@bf=\textbf\def\PY@tc##1{\textcolor[rgb]{0.00,0.50,0.00}{##1}}}
\@namedef{PY@tok@kp}{\def\PY@tc##1{\textcolor[rgb]{0.00,0.50,0.00}{##1}}}
\@namedef{PY@tok@kt}{\def\PY@tc##1{\textcolor[rgb]{0.69,0.00,0.25}{##1}}}
\@namedef{PY@tok@o}{\def\PY@tc##1{\textcolor[rgb]{0.40,0.40,0.40}{##1}}}
\@namedef{PY@tok@ow}{\let\PY@bf=\textbf\def\PY@tc##1{\textcolor[rgb]{0.67,0.13,1.00}{##1}}}
\@namedef{PY@tok@nb}{\def\PY@tc##1{\textcolor[rgb]{0.00,0.50,0.00}{##1}}}
\@namedef{PY@tok@nf}{\def\PY@tc##1{\textcolor[rgb]{0.00,0.00,1.00}{##1}}}
\@namedef{PY@tok@nc}{\let\PY@bf=\textbf\def\PY@tc##1{\textcolor[rgb]{0.00,0.00,1.00}{##1}}}
\@namedef{PY@tok@nn}{\let\PY@bf=\textbf\def\PY@tc##1{\textcolor[rgb]{0.00,0.00,1.00}{##1}}}
\@namedef{PY@tok@ne}{\let\PY@bf=\textbf\def\PY@tc##1{\textcolor[rgb]{0.80,0.25,0.22}{##1}}}
\@namedef{PY@tok@nv}{\def\PY@tc##1{\textcolor[rgb]{0.10,0.09,0.49}{##1}}}
\@namedef{PY@tok@no}{\def\PY@tc##1{\textcolor[rgb]{0.53,0.00,0.00}{##1}}}
\@namedef{PY@tok@nl}{\def\PY@tc##1{\textcolor[rgb]{0.46,0.46,0.00}{##1}}}
\@namedef{PY@tok@ni}{\let\PY@bf=\textbf\def\PY@tc##1{\textcolor[rgb]{0.44,0.44,0.44}{##1}}}
\@namedef{PY@tok@na}{\def\PY@tc##1{\textcolor[rgb]{0.41,0.47,0.13}{##1}}}
\@namedef{PY@tok@nt}{\let\PY@bf=\textbf\def\PY@tc##1{\textcolor[rgb]{0.00,0.50,0.00}{##1}}}
\@namedef{PY@tok@nd}{\def\PY@tc##1{\textcolor[rgb]{0.67,0.13,1.00}{##1}}}
\@namedef{PY@tok@s}{\def\PY@tc##1{\textcolor[rgb]{0.73,0.13,0.13}{##1}}}
\@namedef{PY@tok@sd}{\let\PY@it=\textit\def\PY@tc##1{\textcolor[rgb]{0.73,0.13,0.13}{##1}}}
\@namedef{PY@tok@si}{\let\PY@bf=\textbf\def\PY@tc##1{\textcolor[rgb]{0.64,0.35,0.47}{##1}}}
\@namedef{PY@tok@se}{\let\PY@bf=\textbf\def\PY@tc##1{\textcolor[rgb]{0.67,0.36,0.12}{##1}}}
\@namedef{PY@tok@sr}{\def\PY@tc##1{\textcolor[rgb]{0.64,0.35,0.47}{##1}}}
\@namedef{PY@tok@ss}{\def\PY@tc##1{\textcolor[rgb]{0.10,0.09,0.49}{##1}}}
\@namedef{PY@tok@sx}{\def\PY@tc##1{\textcolor[rgb]{0.00,0.50,0.00}{##1}}}
\@namedef{PY@tok@m}{\def\PY@tc##1{\textcolor[rgb]{0.40,0.40,0.40}{##1}}}
\@namedef{PY@tok@gh}{\let\PY@bf=\textbf\def\PY@tc##1{\textcolor[rgb]{0.00,0.00,0.50}{##1}}}
\@namedef{PY@tok@gu}{\let\PY@bf=\textbf\def\PY@tc##1{\textcolor[rgb]{0.50,0.00,0.50}{##1}}}
\@namedef{PY@tok@gd}{\def\PY@tc##1{\textcolor[rgb]{0.63,0.00,0.00}{##1}}}
\@namedef{PY@tok@gi}{\def\PY@tc##1{\textcolor[rgb]{0.00,0.52,0.00}{##1}}}
\@namedef{PY@tok@gr}{\def\PY@tc##1{\textcolor[rgb]{0.89,0.00,0.00}{##1}}}
\@namedef{PY@tok@ge}{\let\PY@it=\textit}
\@namedef{PY@tok@gs}{\let\PY@bf=\textbf}
\@namedef{PY@tok@gp}{\let\PY@bf=\textbf\def\PY@tc##1{\textcolor[rgb]{0.00,0.00,0.50}{##1}}}
\@namedef{PY@tok@go}{\def\PY@tc##1{\textcolor[rgb]{0.44,0.44,0.44}{##1}}}
\@namedef{PY@tok@gt}{\def\PY@tc##1{\textcolor[rgb]{0.00,0.27,0.87}{##1}}}
\@namedef{PY@tok@err}{\def\PY@bc##1{{\setlength{\fboxsep}{\string -\fboxrule}\fcolorbox[rgb]{1.00,0.00,0.00}{1,1,1}{\strut ##1}}}}
\@namedef{PY@tok@kc}{\let\PY@bf=\textbf\def\PY@tc##1{\textcolor[rgb]{0.00,0.50,0.00}{##1}}}
\@namedef{PY@tok@kd}{\let\PY@bf=\textbf\def\PY@tc##1{\textcolor[rgb]{0.00,0.50,0.00}{##1}}}
\@namedef{PY@tok@kn}{\let\PY@bf=\textbf\def\PY@tc##1{\textcolor[rgb]{0.00,0.50,0.00}{##1}}}
\@namedef{PY@tok@kr}{\let\PY@bf=\textbf\def\PY@tc##1{\textcolor[rgb]{0.00,0.50,0.00}{##1}}}
\@namedef{PY@tok@bp}{\def\PY@tc##1{\textcolor[rgb]{0.00,0.50,0.00}{##1}}}
\@namedef{PY@tok@fm}{\def\PY@tc##1{\textcolor[rgb]{0.00,0.00,1.00}{##1}}}
\@namedef{PY@tok@vc}{\def\PY@tc##1{\textcolor[rgb]{0.10,0.09,0.49}{##1}}}
\@namedef{PY@tok@vg}{\def\PY@tc##1{\textcolor[rgb]{0.10,0.09,0.49}{##1}}}
\@namedef{PY@tok@vi}{\def\PY@tc##1{\textcolor[rgb]{0.10,0.09,0.49}{##1}}}
\@namedef{PY@tok@vm}{\def\PY@tc##1{\textcolor[rgb]{0.10,0.09,0.49}{##1}}}
\@namedef{PY@tok@sa}{\def\PY@tc##1{\textcolor[rgb]{0.73,0.13,0.13}{##1}}}
\@namedef{PY@tok@sb}{\def\PY@tc##1{\textcolor[rgb]{0.73,0.13,0.13}{##1}}}
\@namedef{PY@tok@sc}{\def\PY@tc##1{\textcolor[rgb]{0.73,0.13,0.13}{##1}}}
\@namedef{PY@tok@dl}{\def\PY@tc##1{\textcolor[rgb]{0.73,0.13,0.13}{##1}}}
\@namedef{PY@tok@s2}{\def\PY@tc##1{\textcolor[rgb]{0.73,0.13,0.13}{##1}}}
\@namedef{PY@tok@sh}{\def\PY@tc##1{\textcolor[rgb]{0.73,0.13,0.13}{##1}}}
\@namedef{PY@tok@s1}{\def\PY@tc##1{\textcolor[rgb]{0.73,0.13,0.13}{##1}}}
\@namedef{PY@tok@mb}{\def\PY@tc##1{\textcolor[rgb]{0.40,0.40,0.40}{##1}}}
\@namedef{PY@tok@mf}{\def\PY@tc##1{\textcolor[rgb]{0.40,0.40,0.40}{##1}}}
\@namedef{PY@tok@mh}{\def\PY@tc##1{\textcolor[rgb]{0.40,0.40,0.40}{##1}}}
\@namedef{PY@tok@mi}{\def\PY@tc##1{\textcolor[rgb]{0.40,0.40,0.40}{##1}}}
\@namedef{PY@tok@il}{\def\PY@tc##1{\textcolor[rgb]{0.40,0.40,0.40}{##1}}}
\@namedef{PY@tok@mo}{\def\PY@tc##1{\textcolor[rgb]{0.40,0.40,0.40}{##1}}}
\@namedef{PY@tok@ch}{\let\PY@it=\textit\def\PY@tc##1{\textcolor[rgb]{0.24,0.48,0.48}{##1}}}
\@namedef{PY@tok@cm}{\let\PY@it=\textit\def\PY@tc##1{\textcolor[rgb]{0.24,0.48,0.48}{##1}}}
\@namedef{PY@tok@cpf}{\let\PY@it=\textit\def\PY@tc##1{\textcolor[rgb]{0.24,0.48,0.48}{##1}}}
\@namedef{PY@tok@c1}{\let\PY@it=\textit\def\PY@tc##1{\textcolor[rgb]{0.24,0.48,0.48}{##1}}}
\@namedef{PY@tok@cs}{\let\PY@it=\textit\def\PY@tc##1{\textcolor[rgb]{0.24,0.48,0.48}{##1}}}

\def\PYZbs{\char`\\}
\def\PYZus{\char`\_}
\def\PYZob{\char`\{}
\def\PYZcb{\char`\}}
\def\PYZca{\char`\^}
\def\PYZam{\char`\&}
\def\PYZlt{\char`\<}
\def\PYZgt{\char`\>}
\def\PYZsh{\char`\#}
\def\PYZpc{\char`\%}
\def\PYZdl{\char`\$}
\def\PYZhy{\char`\-}
\def\PYZsq{\char`\'}
\def\PYZdq{\char`\"}
\def\PYZti{\char`\~}
% for compatibility with earlier versions
\def\PYZat{@}
\def\PYZlb{[}
\def\PYZrb{]}
\makeatother


    % For linebreaks inside Verbatim environment from package fancyvrb.
    \makeatletter
        \newbox\Wrappedcontinuationbox
        \newbox\Wrappedvisiblespacebox
        \newcommand*\Wrappedvisiblespace {\textcolor{red}{\textvisiblespace}}
        \newcommand*\Wrappedcontinuationsymbol {\textcolor{red}{\llap{\tiny$\m@th\hookrightarrow$}}}
        \newcommand*\Wrappedcontinuationindent {3ex }
        \newcommand*\Wrappedafterbreak {\kern\Wrappedcontinuationindent\copy\Wrappedcontinuationbox}
        % Take advantage of the already applied Pygments mark-up to insert
        % potential linebreaks for TeX processing.
        %        {, <, #, %, $, ' and ": go to next line.
        %        _, }, ^, &, >, - and ~: stay at end of broken line.
        % Use of \textquotesingle for straight quote.
        \newcommand*\Wrappedbreaksatspecials {%
            \def\PYGZus{\discretionary{\char`\_}{\Wrappedafterbreak}{\char`\_}}%
            \def\PYGZob{\discretionary{}{\Wrappedafterbreak\char`\{}{\char`\{}}%
            \def\PYGZcb{\discretionary{\char`\}}{\Wrappedafterbreak}{\char`\}}}%
            \def\PYGZca{\discretionary{\char`\^}{\Wrappedafterbreak}{\char`\^}}%
            \def\PYGZam{\discretionary{\char`\&}{\Wrappedafterbreak}{\char`\&}}%
            \def\PYGZlt{\discretionary{}{\Wrappedafterbreak\char`\<}{\char`\<}}%
            \def\PYGZgt{\discretionary{\char`\>}{\Wrappedafterbreak}{\char`\>}}%
            \def\PYGZsh{\discretionary{}{\Wrappedafterbreak\char`\#}{\char`\#}}%
            \def\PYGZpc{\discretionary{}{\Wrappedafterbreak\char`\%}{\char`\%}}%
            \def\PYGZdl{\discretionary{}{\Wrappedafterbreak\char`\$}{\char`\$}}%
            \def\PYGZhy{\discretionary{\char`\-}{\Wrappedafterbreak}{\char`\-}}%
            \def\PYGZsq{\discretionary{}{\Wrappedafterbreak\textquotesingle}{\textquotesingle}}%
            \def\PYGZdq{\discretionary{}{\Wrappedafterbreak\char`\"}{\char`\"}}%
            \def\PYGZti{\discretionary{\char`\~}{\Wrappedafterbreak}{\char`\~}}%
        }
        % Some characters . , ; ? ! / are not pygmentized.
        % This macro makes them "active" and they will insert potential linebreaks
        \newcommand*\Wrappedbreaksatpunct {%
            \lccode`\~`\.\lowercase{\def~}{\discretionary{\hbox{\char`\.}}{\Wrappedafterbreak}{\hbox{\char`\.}}}%
            \lccode`\~`\,\lowercase{\def~}{\discretionary{\hbox{\char`\,}}{\Wrappedafterbreak}{\hbox{\char`\,}}}%
            \lccode`\~`\;\lowercase{\def~}{\discretionary{\hbox{\char`\;}}{\Wrappedafterbreak}{\hbox{\char`\;}}}%
            \lccode`\~`\:\lowercase{\def~}{\discretionary{\hbox{\char`\:}}{\Wrappedafterbreak}{\hbox{\char`\:}}}%
            \lccode`\~`\?\lowercase{\def~}{\discretionary{\hbox{\char`\?}}{\Wrappedafterbreak}{\hbox{\char`\?}}}%
            \lccode`\~`\!\lowercase{\def~}{\discretionary{\hbox{\char`\!}}{\Wrappedafterbreak}{\hbox{\char`\!}}}%
            \lccode`\~`\/\lowercase{\def~}{\discretionary{\hbox{\char`\/}}{\Wrappedafterbreak}{\hbox{\char`\/}}}%
            \catcode`\.\active
            \catcode`\,\active
            \catcode`\;\active
            \catcode`\:\active
            \catcode`\?\active
            \catcode`\!\active
            \catcode`\/\active
            \lccode`\~`\~
        }
    \makeatother

    \let\OriginalVerbatim=\Verbatim
    \makeatletter
    \renewcommand{\Verbatim}[1][1]{%
        %\parskip\z@skip
        \sbox\Wrappedcontinuationbox {\Wrappedcontinuationsymbol}%
        \sbox\Wrappedvisiblespacebox {\FV@SetupFont\Wrappedvisiblespace}%
        \def\FancyVerbFormatLine ##1{\hsize\linewidth
            \vtop{\raggedright\hyphenpenalty\z@\exhyphenpenalty\z@
                \doublehyphendemerits\z@\finalhyphendemerits\z@
                \strut ##1\strut}%
        }%
        % If the linebreak is at a space, the latter will be displayed as visible
        % space at end of first line, and a continuation symbol starts next line.
        % Stretch/shrink are however usually zero for typewriter font.
        \def\FV@Space {%
            \nobreak\hskip\z@ plus\fontdimen3\font minus\fontdimen4\font
            \discretionary{\copy\Wrappedvisiblespacebox}{\Wrappedafterbreak}
            {\kern\fontdimen2\font}%
        }%

        % Allow breaks at special characters using \PYG... macros.
        \Wrappedbreaksatspecials
        % Breaks at punctuation characters . , ; ? ! and / need catcode=\active
        \OriginalVerbatim[#1,codes*=\Wrappedbreaksatpunct]%
    }
    \makeatother

    % Exact colors from NB
    \definecolor{incolor}{HTML}{303F9F}
    \definecolor{outcolor}{HTML}{D84315}
    \definecolor{cellborder}{HTML}{CFCFCF}
    \definecolor{cellbackground}{HTML}{F7F7F7}

    % prompt
    \makeatletter
    \newcommand{\boxspacing}{\kern\kvtcb@left@rule\kern\kvtcb@boxsep}
    \makeatother
    \newcommand{\prompt}[4]{
        {\ttfamily\llap{{\color{#2}[#3]:\hspace{3pt}#4}}\vspace{-\baselineskip}}
    }
    

    
    % Prevent overflowing lines due to hard-to-break entities
    \sloppy
    % Setup hyperref package
    \hypersetup{
      breaklinks=true,  % so long urls are correctly broken across lines
      colorlinks=true,
      urlcolor=urlcolor,
      linkcolor=linkcolor,
      citecolor=citecolor,
      }
    % Slightly bigger margins than the latex defaults
    
    \geometry{verbose,tmargin=1in,bmargin=1in,lmargin=1in,rmargin=1in}
    
    

\begin{document}
    
    \maketitle
    
    

    
    \hypertarget{spectroscopic-data-reduction-part-2-wavelength-calibration}{%
\section{Spectroscopic Data Reduction Part 2: Wavelength
Calibration}\label{spectroscopic-data-reduction-part-2-wavelength-calibration}}

This notebook assumes you've completed the Spectroscopic Trace process
(see \href{1-SpectroscopicTraceTutorial.ipynb}{Part 1}) and have a trace
model handy.

\hypertarget{reducciuxf3n-de-datos-espectroscuxf3picos-parte-2-calibraciuxf3n-en-longitud-de-onda}{%
\section{Reducción de Datos Espectroscópicos Parte 2: Calibración en
Longitud de
Onda}\label{reducciuxf3n-de-datos-espectroscuxf3picos-parte-2-calibraciuxf3n-en-longitud-de-onda}}

\textbf{Resumen:} En este cuaderno continuamos con la calibración
espectral (wavelength calibration). Vamos a: - cargar imágenes de
espectros y convertirlas a vectores de intensidad vs.~píxel, -
identificar picos (líneas espectrales) en las imágenes de referencia
(He, Ne, Kr, Hg, \ldots), - buscar las longitudes de onda conocidas
asociadas a esos picos (usando catálogos o la base de datos NIST), -
ajustar un polinomio que convierta píxeles → longitud de onda, - aplicar
esa calibración a espectros observacionales para poder medir
características (líneas, anchos, etc.).

\textbf{Autores y referencia del tutorial original:} Adam Ginsburg,
Kelle Cruz, Lia Corrales, Jonathan Sick, Adrian Price-Whelan

    \textbf{Requerimientos (traducción y explicación):}

Se lee \texttt{requirements.txt} y muestra los paquetes necesarios para
ejecutar el cuaderno. - Asegúrate de tener instalados los paquetes
listados (por ejemplo: astropy, numpy, matplotlib, scipy, photutils,
specutils, PIL/Pillow). - Si falta algún paquete, instálalo con
\texttt{pip\ install\ -r\ requirements.txt} o equivalente en tu entorno.
- Si trabajas en un entorno virtual, confirma que Jupyter/Notebook esté
usando ese mismo entorno (kernel).

    \begin{tcolorbox}[breakable, size=fbox, boxrule=1pt, pad at break*=1mm,colback=cellbackground, colframe=cellborder]
\prompt{In}{incolor}{1}{\boxspacing}
\begin{Verbatim}[commandchars=\\\{\}]
\PY{k}{with} \PY{n+nb}{open}\PY{p}{(}\PY{l+s+s1}{\PYZsq{}}\PY{l+s+s1}{requirements.txt}\PY{l+s+s1}{\PYZsq{}}\PY{p}{)} \PY{k}{as} \PY{n}{f}\PY{p}{:}
    \PY{n+nb}{print}\PY{p}{(}\PY{l+s+sa}{f}\PY{l+s+s2}{\PYZdq{}}\PY{l+s+s2}{Required packages for this notebook:}\PY{l+s+se}{\PYZbs{}n}\PY{l+s+si}{\PYZob{}}\PY{n}{f}\PY{o}{.}\PY{n}{read}\PY{p}{(}\PY{p}{)}\PY{l+s+si}{\PYZcb{}}\PY{l+s+s2}{\PYZdq{}}\PY{p}{)}
\end{Verbatim}
\end{tcolorbox}

    \begin{Verbatim}[commandchars=\\\{\}]
Required packages for this notebook:
astropy
astroquery>=0.4.8.dev9474  \# 2024-09-24 pinned for Gaia column capitalization
issue
IPython
numpy
pillow
matplotlib

    \end{Verbatim}

    \textbf{Imports --- explicación paso a paso y traducción:}

La celda de imports carga librerías necesarias. Resumen de las más
importantes:

\begin{itemize}
\tightlist
\item
  \texttt{PIL.Image} (Pillow): para abrir imágenes de espectros (jpeg,
  png).
\item
  \texttt{numpy}: manejo de arrays y cálculos numéricos.
\item
  \texttt{matplotlib.pyplot}: para graficar espectros, imágenes,
  perfiles.
\item
  \texttt{astropy.units} y \texttt{astropy.modeling.polynomial}:
  herramientas astronómicas y modelos polinómicos (útiles para la
  calibración).
\item
  \texttt{scipy.signal} (p.~ej. \texttt{find\_peaks}): para detectar
  picos en el espectro.
\item
  \texttt{photutils} o \texttt{specutils} si está presente: rutinas
  específicas para fotometría/apertura/line finding.
\end{itemize}

    \begin{tcolorbox}[breakable, size=fbox, boxrule=1pt, pad at break*=1mm,colback=cellbackground, colframe=cellborder]
\prompt{In}{incolor}{2}{\boxspacing}
\begin{Verbatim}[commandchars=\\\{\}]
\PY{k+kn}{from} \PY{n+nn}{PIL} \PY{k+kn}{import} \PY{n}{Image} \PY{k}{as} \PY{n}{PILImage}
\PY{k+kn}{import} \PY{n+nn}{numpy} \PY{k}{as} \PY{n+nn}{np}
\PY{k+kn}{import} \PY{n+nn}{matplotlib}\PY{n+nn}{.}\PY{n+nn}{pyplot} \PY{k}{as} \PY{n+nn}{plt}
\PY{n}{plt}\PY{o}{.}\PY{n}{rcParams}\PY{p}{[}\PY{l+s+s1}{\PYZsq{}}\PY{l+s+s1}{image.origin}\PY{l+s+s1}{\PYZsq{}}\PY{p}{]} \PY{o}{=} \PY{l+s+s1}{\PYZsq{}}\PY{l+s+s1}{lower}\PY{l+s+s1}{\PYZsq{}}  \PY{c+c1}{\PYZsh{} we want to show images, not matrices, so we set the origin to the lower\PYZhy{}left}
\PY{n}{plt}\PY{o}{.}\PY{n}{style}\PY{o}{.}\PY{n}{use}\PY{p}{(}\PY{l+s+s1}{\PYZsq{}}\PY{l+s+s1}{dark\PYZus{}background}\PY{l+s+s1}{\PYZsq{}}\PY{p}{)}  \PY{c+c1}{\PYZsh{} Optional configuration: if run, this will look nice on dark background notebooks}
\end{Verbatim}
\end{tcolorbox}

    \begin{tcolorbox}[breakable, size=fbox, boxrule=1pt, pad at break*=1mm,colback=cellbackground, colframe=cellborder]
\prompt{In}{incolor}{3}{\boxspacing}
\begin{Verbatim}[commandchars=\\\{\}]
\PY{k+kn}{from} \PY{n+nn}{astropy} \PY{k+kn}{import} \PY{n}{units} \PY{k}{as} \PY{n}{u}
\PY{k+kn}{from} \PY{n+nn}{astropy}\PY{n+nn}{.}\PY{n+nn}{modeling}\PY{n+nn}{.}\PY{n+nn}{polynomial} \PY{k+kn}{import} \PY{n}{Polynomial1D}
\PY{k+kn}{from} \PY{n+nn}{astropy}\PY{n+nn}{.}\PY{n+nn}{modeling}\PY{n+nn}{.}\PY{n+nn}{models} \PY{k+kn}{import} \PY{n}{Gaussian1D}\PY{p}{,} \PY{n}{Linear1D}
\PY{k+kn}{from} \PY{n+nn}{astropy}\PY{n+nn}{.}\PY{n+nn}{modeling}\PY{n+nn}{.}\PY{n+nn}{fitting} \PY{k+kn}{import} \PY{n}{LinearLSQFitter}
\PY{k+kn}{from} \PY{n+nn}{IPython}\PY{n+nn}{.}\PY{n+nn}{display} \PY{k+kn}{import} \PY{n}{Image}
\PY{c+c1}{\PYZsh{} astroquery provides an interface to the NIST atomic line database}
\PY{k+kn}{from} \PY{n+nn}{astroquery}\PY{n+nn}{.}\PY{n+nn}{nist} \PY{k+kn}{import} \PY{n}{Nist}
\end{Verbatim}
\end{tcolorbox}

    \begin{tcolorbox}[breakable, size=fbox, boxrule=1pt, pad at break*=1mm,colback=cellbackground, colframe=cellborder]
\prompt{In}{incolor}{ }{\boxspacing}
\begin{Verbatim}[commandchars=\\\{\}]
\PY{c+c1}{\PYZsh{} Cell 1: Imports y configuración}
\PY{k+kn}{import} \PY{n+nn}{os}
\PY{k+kn}{from} \PY{n+nn}{pathlib} \PY{k+kn}{import} \PY{n}{Path}
\PY{k+kn}{import} \PY{n+nn}{numpy} \PY{k}{as} \PY{n+nn}{np}
\PY{k+kn}{import} \PY{n+nn}{matplotlib}\PY{n+nn}{.}\PY{n+nn}{pyplot} \PY{k}{as} \PY{n+nn}{plt}
\PY{k+kn}{from} \PY{n+nn}{PIL} \PY{k+kn}{import} \PY{n}{Image}   
\PY{k+kn}{from} \PY{n+nn}{scipy}\PY{n+nn}{.}\PY{n+nn}{signal} \PY{k+kn}{import} \PY{n}{find\PYZus{}peaks}\PY{p}{,} \PY{n}{savgol\PYZus{}filter}
\PY{k+kn}{from} \PY{n+nn}{scipy}\PY{n+nn}{.}\PY{n+nn}{optimize} \PY{k+kn}{import} \PY{n}{curve\PYZus{}fit}
\PY{k+kn}{import} \PY{n+nn}{pandas} \PY{k}{as} \PY{n+nn}{pd}

\PY{c+c1}{\PYZsh{} Mostrar figuras inline en Jupyter}
\PY{o}{\PYZpc{}}\PY{k}{matplotlib} inline
\PY{n}{plt}\PY{o}{.}\PY{n}{rcParams}\PY{p}{[}\PY{l+s+s1}{\PYZsq{}}\PY{l+s+s1}{figure.figsize}\PY{l+s+s1}{\PYZsq{}}\PY{p}{]} \PY{o}{=} \PY{p}{(}\PY{l+m+mi}{10}\PY{p}{,}\PY{l+m+mi}{4}\PY{p}{)}
\end{Verbatim}
\end{tcolorbox}

    \textbf{Explicación de la visualización de la imagen:}

\begin{itemize}
\tightlist
\item
  Se carga la imagen (por ejemplo: imagen del espectro de Ne, He, Kr).
\item
  Se muestra con \texttt{plt.imshow(...)} para inspeccionar visualmente
  la distribución espacial de la luz.
\item
  Objetivo: localizar la franja donde está el espectro (usualmente una
  franja horizontal/vertical), verificar rotación, saturación o
  artefactos.
\end{itemize}

\textbf{Paso a paso:} 1. Mira la orientación (¿el espectro está
horizontal o vertical?). Si no está alineado, podría ser necesario rotar
o recortar la imagen. 2. Localiza la franja donde la intensidad varía
(esa es la dispersión). 3. Si hay líneas verticales brillantes y fondo
estable: es buena señal para extracción.

    \begin{tcolorbox}[breakable, size=fbox, boxrule=1pt, pad at break*=1mm,colback=cellbackground, colframe=cellborder]
\prompt{In}{incolor}{4}{\boxspacing}
\begin{Verbatim}[commandchars=\\\{\}]
\PY{n}{hg\PYZus{}filename} \PY{o}{=} \PY{l+s+s2}{\PYZdq{}}\PY{l+s+s2}{Mercurio.bmp}\PY{l+s+s2}{\PYZdq{}}
\PY{n}{kr\PYZus{}filename} \PY{o}{=} \PY{l+s+s2}{\PYZdq{}}\PY{l+s+s2}{Kripton.bmp}\PY{l+s+s2}{\PYZdq{}}
\PY{n}{ne\PYZus{}filename} \PY{o}{=} \PY{l+s+s2}{\PYZdq{}}\PY{l+s+s2}{Neon.bmp}\PY{l+s+s2}{\PYZdq{}}
\PY{n}{he\PYZus{}filename} \PY{o}{=} \PY{l+s+s2}{\PYZdq{}}\PY{l+s+s2}{Helio.bmp}\PY{l+s+s2}{\PYZdq{}}
\end{Verbatim}
\end{tcolorbox}

    img = PILImage.open(`file.jpg') → abre el archivo de imagen.

arr = np.array(img) → convierte la imagen a un array NumPy (shape: filas
x columnas x canales).

plt.imshow(arr, origin=`lower') → muestra la imagen para inspección
visual.

\textbf{Explicación paso a paso:}

Abre la imagen del espectro (puede ser foto del foco del espectrógrafo o
captura del detector).

Convierte la imagen a array para poder manipular pixeles numéricamente.

Muestra la imagen para verificar orientación, saturación, y la presencia
de la franja espectral.

    \begin{tcolorbox}[breakable, size=fbox, boxrule=1pt, pad at break*=1mm,colback=cellbackground, colframe=cellborder]
\prompt{In}{incolor}{5}{\boxspacing}
\begin{Verbatim}[commandchars=\\\{\}]
\PY{n}{hg\PYZus{}image} \PY{o}{=} \PY{n}{np}\PY{o}{.}\PY{n}{array}\PY{p}{(}\PY{n}{PILImage}\PY{o}{.}\PY{n}{open}\PY{p}{(}\PY{n}{hg\PYZus{}filename}\PY{p}{)}\PY{p}{)}
\PY{n}{kr\PYZus{}image} \PY{o}{=} \PY{n}{np}\PY{o}{.}\PY{n}{array}\PY{p}{(}\PY{n}{PILImage}\PY{o}{.}\PY{n}{open}\PY{p}{(}\PY{n}{kr\PYZus{}filename}\PY{p}{)}\PY{p}{)}
\PY{n}{ne\PYZus{}image} \PY{o}{=} \PY{n}{np}\PY{o}{.}\PY{n}{array}\PY{p}{(}\PY{n}{PILImage}\PY{o}{.}\PY{n}{open}\PY{p}{(}\PY{n}{ne\PYZus{}filename}\PY{p}{)}\PY{p}{)}
\PY{n}{he\PYZus{}image} \PY{o}{=} \PY{n}{np}\PY{o}{.}\PY{n}{array}\PY{p}{(}\PY{n}{PILImage}\PY{o}{.}\PY{n}{open}\PY{p}{(}\PY{n}{he\PYZus{}filename}\PY{p}{)}\PY{p}{)}
\end{Verbatim}
\end{tcolorbox}

    \textbf{Convertir JPEG a BMP}

Se convierten las imágenes ``He.jpeg'', ``Kr.jpeg'', ``Hg.jpeg'',
``Ne.jpeg'' al formato bmp

    \begin{tcolorbox}[breakable, size=fbox, boxrule=1pt, pad at break*=1mm,colback=cellbackground, colframe=cellborder]
\prompt{In}{incolor}{ }{\boxspacing}
\begin{Verbatim}[commandchars=\\\{\}]
\PY{c+c1}{\PYZsh{} Cell 2: Convertir JPEG a BMP (opcional)}
\PY{n}{jpg\PYZus{}files} \PY{o}{=} \PY{p}{[}\PY{l+s+s2}{\PYZdq{}}\PY{l+s+s2}{He.jpeg}\PY{l+s+s2}{\PYZdq{}}\PY{p}{,} \PY{l+s+s2}{\PYZdq{}}\PY{l+s+s2}{Kr.jpeg}\PY{l+s+s2}{\PYZdq{}}\PY{p}{,} \PY{l+s+s2}{\PYZdq{}}\PY{l+s+s2}{Hg.jpeg}\PY{l+s+s2}{\PYZdq{}}\PY{p}{,} \PY{l+s+s2}{\PYZdq{}}\PY{l+s+s2}{Ne.jpeg}\PY{l+s+s2}{\PYZdq{}}\PY{p}{]}
\PY{n}{out\PYZus{}format} \PY{o}{=} \PY{l+s+s2}{\PYZdq{}}\PY{l+s+s2}{bmp}\PY{l+s+s2}{\PYZdq{}}  \PY{c+c1}{\PYZsh{} o \PYZsq{}png\PYZsq{}}

\PY{k}{for} \PY{n}{fname} \PY{o+ow}{in} \PY{n}{jpg\PYZus{}files}\PY{p}{:}
    \PY{k}{if} \PY{o+ow}{not} \PY{n}{os}\PY{o}{.}\PY{n}{path}\PY{o}{.}\PY{n}{exists}\PY{p}{(}\PY{n}{fname}\PY{p}{)}\PY{p}{:}
        \PY{n+nb}{print}\PY{p}{(}\PY{l+s+sa}{f}\PY{l+s+s2}{\PYZdq{}}\PY{l+s+s2}{Archivo no encontrado: }\PY{l+s+si}{\PYZob{}}\PY{n}{fname}\PY{l+s+si}{\PYZcb{}}\PY{l+s+s2}{\PYZdq{}}\PY{p}{)}
        \PY{k}{continue}
    \PY{n}{base} \PY{o}{=} \PY{n}{Path}\PY{p}{(}\PY{n}{fname}\PY{p}{)}\PY{o}{.}\PY{n}{stem}
    \PY{n}{outname} \PY{o}{=} \PY{l+s+sa}{f}\PY{l+s+s2}{\PYZdq{}}\PY{l+s+si}{\PYZob{}}\PY{n}{base}\PY{l+s+si}{\PYZcb{}}\PY{l+s+s2}{.}\PY{l+s+si}{\PYZob{}}\PY{n}{out\PYZus{}format}\PY{l+s+si}{\PYZcb{}}\PY{l+s+s2}{\PYZdq{}}
    \PY{k}{if} \PY{o+ow}{not} \PY{n}{os}\PY{o}{.}\PY{n}{path}\PY{o}{.}\PY{n}{exists}\PY{p}{(}\PY{n}{outname}\PY{p}{)}\PY{p}{:}
        \PY{n}{img} \PY{o}{=} \PY{n}{Image}\PY{o}{.}\PY{n}{open}\PY{p}{(}\PY{n}{fname}\PY{p}{)}\PY{o}{.}\PY{n}{convert}\PY{p}{(}\PY{l+s+s1}{\PYZsq{}}\PY{l+s+s1}{RGB}\PY{l+s+s1}{\PYZsq{}}\PY{p}{)}
        \PY{n}{img}\PY{o}{.}\PY{n}{save}\PY{p}{(}\PY{n}{outname}\PY{p}{)}
        \PY{n+nb}{print}\PY{p}{(}\PY{l+s+sa}{f}\PY{l+s+s2}{\PYZdq{}}\PY{l+s+s2}{Guardado: }\PY{l+s+si}{\PYZob{}}\PY{n}{outname}\PY{l+s+si}{\PYZcb{}}\PY{l+s+s2}{\PYZdq{}}\PY{p}{)}
    \PY{k}{else}\PY{p}{:}
        \PY{n+nb}{print}\PY{p}{(}\PY{l+s+sa}{f}\PY{l+s+s2}{\PYZdq{}}\PY{l+s+s2}{Ya existe: }\PY{l+s+si}{\PYZob{}}\PY{n}{outname}\PY{l+s+si}{\PYZcb{}}\PY{l+s+s2}{\PYZdq{}}\PY{p}{)}
\end{Verbatim}
\end{tcolorbox}

    \textbf{Se carga la imagen y se convierte a 1D:}

Como las imágenes poseen varios colores se convierten a arreglos
unidimensionales para poder procesarlas

    \begin{tcolorbox}[breakable, size=fbox, boxrule=1pt, pad at break*=1mm,colback=cellbackground, colframe=cellborder]
\prompt{In}{incolor}{6}{\boxspacing}
\begin{Verbatim}[commandchars=\\\{\}]
\PY{c+c1}{\PYZsh{} Cell 3: Funciones auxiliares para cargar imagen y extraer perfil espectral}
\PY{k}{def} \PY{n+nf}{load\PYZus{}image\PYZus{}gray}\PY{p}{(}\PY{n}{path}\PY{p}{)}\PY{p}{:}
    \PY{n}{img} \PY{o}{=} \PY{n}{Image}\PY{o}{.}\PY{n}{open}\PY{p}{(}\PY{n}{path}\PY{p}{)}\PY{o}{.}\PY{n}{convert}\PY{p}{(}\PY{l+s+s1}{\PYZsq{}}\PY{l+s+s1}{L}\PY{l+s+s1}{\PYZsq{}}\PY{p}{)}
    \PY{n}{arr} \PY{o}{=} \PY{n}{np}\PY{o}{.}\PY{n}{array}\PY{p}{(}\PY{n}{img}\PY{p}{)}\PY{o}{.}\PY{n}{astype}\PY{p}{(}\PY{n+nb}{float}\PY{p}{)}
    \PY{n}{arr} \PY{o}{=} \PY{p}{(}\PY{n}{arr} \PY{o}{\PYZhy{}} \PY{n}{arr}\PY{o}{.}\PY{n}{min}\PY{p}{(}\PY{p}{)}\PY{p}{)} \PY{o}{/} \PY{p}{(}\PY{n}{arr}\PY{o}{.}\PY{n}{max}\PY{p}{(}\PY{p}{)} \PY{o}{\PYZhy{}} \PY{n}{arr}\PY{o}{.}\PY{n}{min}\PY{p}{(}\PY{p}{)} \PY{o}{+} \PY{l+m+mf}{1e\PYZhy{}12}\PY{p}{)}
    \PY{k}{return} \PY{n}{arr}

\PY{k}{def} \PY{n+nf}{spectral\PYZus{}profile}\PY{p}{(}\PY{n}{img\PYZus{}arr}\PY{p}{,} \PY{n}{axis}\PY{o}{=}\PY{l+m+mi}{0}\PY{p}{,} \PY{n}{crop}\PY{o}{=}\PY{k+kc}{None}\PY{p}{)}\PY{p}{:}
    \PY{k}{if} \PY{n}{crop} \PY{o+ow}{is} \PY{o+ow}{not} \PY{k+kc}{None}\PY{p}{:}
        \PY{n}{r0}\PY{p}{,}\PY{n}{r1}\PY{p}{,}\PY{n}{c0}\PY{p}{,}\PY{n}{c1} \PY{o}{=} \PY{n}{crop}
        \PY{n}{img\PYZus{}arr} \PY{o}{=} \PY{n}{img\PYZus{}arr}\PY{p}{[}\PY{n}{r0}\PY{p}{:}\PY{n}{r1}\PY{p}{,} \PY{n}{c0}\PY{p}{:}\PY{n}{c1}\PY{p}{]}
    \PY{k}{return} \PY{n}{img\PYZus{}arr}\PY{o}{.}\PY{n}{sum}\PY{p}{(}\PY{n}{axis}\PY{o}{=}\PY{n}{axis}\PY{p}{)}

\PY{k}{def} \PY{n+nf}{smooth\PYZus{}profile}\PY{p}{(}\PY{n}{profile}\PY{p}{,} \PY{n}{window}\PY{o}{=}\PY{l+m+mi}{31}\PY{p}{,} \PY{n}{poly}\PY{o}{=}\PY{l+m+mi}{3}\PY{p}{)}\PY{p}{:}
    \PY{k}{if} \PY{n+nb}{len}\PY{p}{(}\PY{n}{profile}\PY{p}{)} \PY{o}{\PYZlt{}} \PY{n}{window}\PY{p}{:}
        \PY{n}{window} \PY{o}{=} \PY{n+nb}{max}\PY{p}{(}\PY{l+m+mi}{3}\PY{p}{,} \PY{p}{(}\PY{n+nb}{len}\PY{p}{(}\PY{n}{profile}\PY{p}{)}\PY{o}{/}\PY{o}{/}\PY{l+m+mi}{2}\PY{p}{)}\PY{o}{*}\PY{l+m+mi}{2}\PY{o}{+}\PY{l+m+mi}{1}\PY{p}{)}
    \PY{k}{return} \PY{n}{savgol\PYZus{}filter}\PY{p}{(}\PY{n}{profile}\PY{p}{,} \PY{n}{window}\PY{p}{,} \PY{n}{poly}\PY{p}{)}
\end{Verbatim}
\end{tcolorbox}

    \begin{tcolorbox}[breakable, size=fbox, boxrule=1pt, pad at break*=1mm,colback=cellbackground, colframe=cellborder]
\prompt{In}{incolor}{8}{\boxspacing}
\begin{Verbatim}[commandchars=\\\{\}]
\PY{n}{hg\PYZus{}gray} \PY{o}{=} \PY{n}{load\PYZus{}image\PYZus{}gray}\PY{p}{(}\PY{n}{hg\PYZus{}filename}\PY{p}{)}
\PY{n}{kr\PYZus{}gray} \PY{o}{=} \PY{n}{load\PYZus{}image\PYZus{}gray}\PY{p}{(}\PY{n}{kr\PYZus{}filename}\PY{p}{)}
\PY{n}{ne\PYZus{}gray} \PY{o}{=} \PY{n}{load\PYZus{}image\PYZus{}gray}\PY{p}{(}\PY{n}{ne\PYZus{}filename}\PY{p}{)}
\PY{n}{he\PYZus{}gray} \PY{o}{=} \PY{n}{load\PYZus{}image\PYZus{}gray}\PY{p}{(}\PY{n}{he\PYZus{}filename}\PY{p}{)}
\end{Verbatim}
\end{tcolorbox}

    \textbf{Impresión de las imágenes unidimensionales}

    \begin{tcolorbox}[breakable, size=fbox, boxrule=1pt, pad at break*=1mm,colback=cellbackground, colframe=cellborder]
\prompt{In}{incolor}{9}{\boxspacing}
\begin{Verbatim}[commandchars=\\\{\}]
\PY{n}{plt}\PY{o}{.}\PY{n}{imshow}\PY{p}{(}\PY{n}{hg\PYZus{}gray}\PY{p}{)}\PY{p}{;}
\end{Verbatim}
\end{tcolorbox}

    \begin{center}
    \adjustimage{max size={0.9\linewidth}{0.9\paperheight}}{output_17_0.png}
    \end{center}
    { \hspace*{\fill} \\}
    
    \begin{tcolorbox}[breakable, size=fbox, boxrule=1pt, pad at break*=1mm,colback=cellbackground, colframe=cellborder]
\prompt{In}{incolor}{10}{\boxspacing}
\begin{Verbatim}[commandchars=\\\{\}]
\PY{n}{plt}\PY{o}{.}\PY{n}{imshow}\PY{p}{(}\PY{n}{kr\PYZus{}gray}\PY{p}{)}\PY{p}{;}
\end{Verbatim}
\end{tcolorbox}

    \begin{center}
    \adjustimage{max size={0.9\linewidth}{0.9\paperheight}}{output_18_0.png}
    \end{center}
    { \hspace*{\fill} \\}
    
    \begin{tcolorbox}[breakable, size=fbox, boxrule=1pt, pad at break*=1mm,colback=cellbackground, colframe=cellborder]
\prompt{In}{incolor}{11}{\boxspacing}
\begin{Verbatim}[commandchars=\\\{\}]
\PY{n}{plt}\PY{o}{.}\PY{n}{imshow}\PY{p}{(}\PY{n}{ne\PYZus{}gray}\PY{p}{)}\PY{p}{;}
\end{Verbatim}
\end{tcolorbox}

    \begin{center}
    \adjustimage{max size={0.9\linewidth}{0.9\paperheight}}{output_19_0.png}
    \end{center}
    { \hspace*{\fill} \\}
    
    \begin{tcolorbox}[breakable, size=fbox, boxrule=1pt, pad at break*=1mm,colback=cellbackground, colframe=cellborder]
\prompt{In}{incolor}{12}{\boxspacing}
\begin{Verbatim}[commandchars=\\\{\}]
\PY{n}{plt}\PY{o}{.}\PY{n}{imshow}\PY{p}{(}\PY{n}{he\PYZus{}gray}\PY{p}{)}
\end{Verbatim}
\end{tcolorbox}

            \begin{tcolorbox}[breakable, size=fbox, boxrule=.5pt, pad at break*=1mm, opacityfill=0]
\prompt{Out}{outcolor}{12}{\boxspacing}
\begin{Verbatim}[commandchars=\\\{\}]
<matplotlib.image.AxesImage at 0x7f2f39e0fc10>
\end{Verbatim}
\end{tcolorbox}
        
    \begin{center}
    \adjustimage{max size={0.9\linewidth}{0.9\paperheight}}{output_20_1.png}
    \end{center}
    { \hspace*{\fill} \\}
    
    Podemos observar únicamente una línea espectral haciendo zoom

    \begin{tcolorbox}[breakable, size=fbox, boxrule=1pt, pad at break*=1mm,colback=cellbackground, colframe=cellborder]
\prompt{In}{incolor}{13}{\boxspacing}
\begin{Verbatim}[commandchars=\\\{\}]
\PY{n}{plt}\PY{o}{.}\PY{n}{imshow}\PY{p}{(}\PY{n}{hg\PYZus{}gray}\PY{p}{[}\PY{l+m+mi}{100}\PY{p}{:}\PY{l+m+mi}{300}\PY{p}{,}\PY{l+m+mi}{280}\PY{p}{:}\PY{l+m+mi}{340}\PY{p}{]}\PY{p}{)}\PY{p}{;}
\end{Verbatim}
\end{tcolorbox}

    \begin{center}
    \adjustimage{max size={0.9\linewidth}{0.9\paperheight}}{output_22_0.png}
    \end{center}
    { \hspace*{\fill} \\}
    
    Nosotros recreamos nuestro modelo de traza del
\href{Spectroscopic\%20Trace\%20Tutorial.ipynb}{Tutorial de Trazado
Espectroscópico} usando los modelos ajustados.

Podríamos haber usado la traza empírica directamente, lo cual podría
resultar en características de ruido ligeramente mejoradas, pero por
simplicidad ---y para hacer que ambos cuadernos sean utilizables de
manera independiente--- usamos aquí los modelos polinómicos y Gaussianos
ajustados.

    \begin{tcolorbox}[breakable, size=fbox, boxrule=1pt, pad at break*=1mm,colback=cellbackground, colframe=cellborder]
\prompt{In}{incolor}{89}{\boxspacing}
\begin{Verbatim}[commandchars=\\\{\}]
\PY{n}{trace\PYZus{}model} \PY{o}{=} \PY{n}{Polynomial1D}\PY{p}{(}\PY{n}{degree}\PY{o}{=}\PY{l+m+mi}{3}\PY{p}{,} \PY{n}{c0}\PY{o}{=}\PY{l+m+mf}{276.03448116}\PY{p}{,} \PY{n}{c1}\PY{o}{=}\PY{l+m+mf}{0.0091462}\PY{p}{,} \PY{n}{c2}\PY{o}{=}\PY{l+m+mf}{0.00002231}\PY{p}{,} \PY{n}{c3}\PY{o}{=}\PY{o}{\PYZhy{}}\PY{l+m+mf}{0.00000001}\PY{p}{)}
\PY{n}{trace\PYZus{}profile\PYZus{}model} \PY{o}{=} \PY{n}{Gaussian1D}\PY{p}{(}\PY{n}{amplitude}\PY{o}{=}\PY{l+m+mf}{123.84846797}\PY{p}{,} \PY{n}{mean}\PY{o}{=}\PY{l+m+mf}{0.17719819}\PY{p}{,} \PY{n}{stddev}\PY{o}{=}\PY{l+m+mf}{5.10872134}\PY{p}{)}
\PY{n}{xaxis} \PY{o}{=} \PY{n}{np}\PY{o}{.}\PY{n}{arange}\PY{p}{(}\PY{n}{he\PYZus{}gray}\PY{o}{.}\PY{n}{shape}\PY{p}{[}\PY{l+m+mi}{1}\PY{p}{]}\PY{p}{)}
\PY{n}{trace\PYZus{}center} \PY{o}{=} \PY{n}{trace\PYZus{}model}\PY{p}{(}\PY{n}{xaxis}\PY{p}{)}
\PY{n}{npixels\PYZus{}to\PYZus{}cut}\PY{o}{=}\PY{l+m+mi}{15}
\PY{n}{yaxis} \PY{o}{=} \PY{n}{np}\PY{o}{.}\PY{n}{arange}\PY{p}{(}\PY{o}{\PYZhy{}}\PY{n}{npixels\PYZus{}to\PYZus{}cut}\PY{p}{,} \PY{n}{npixels\PYZus{}to\PYZus{}cut}\PY{p}{)}
\PY{n}{model\PYZus{}trace\PYZus{}profile} \PY{o}{=} \PY{n}{trace\PYZus{}profile\PYZus{}model}\PY{p}{(}\PY{n}{yaxis}\PY{p}{)}
\end{Verbatim}
\end{tcolorbox}

    \begin{tcolorbox}[breakable, size=fbox, boxrule=1pt, pad at break*=1mm,colback=cellbackground, colframe=cellborder]
\prompt{In}{incolor}{90}{\boxspacing}
\begin{Verbatim}[commandchars=\\\{\}]
\PY{n}{xaxis\PYZus{}hg} \PY{o}{=} \PY{n}{np}\PY{o}{.}\PY{n}{arange}\PY{p}{(}\PY{n}{hg\PYZus{}gray}\PY{o}{.}\PY{n}{shape}\PY{p}{[}\PY{l+m+mi}{1}\PY{p}{]}\PY{p}{)}
\PY{n}{trace\PYZus{}center\PYZus{}hg} \PY{o}{=} \PY{n}{trace\PYZus{}model}\PY{p}{(}\PY{n}{xaxis\PYZus{}hg}\PY{p}{)}
\end{Verbatim}
\end{tcolorbox}

    \begin{tcolorbox}[breakable, size=fbox, boxrule=1pt, pad at break*=1mm,colback=cellbackground, colframe=cellborder]
\prompt{In}{incolor}{91}{\boxspacing}
\begin{Verbatim}[commandchars=\\\{\}]
\PY{n}{xaxis\PYZus{}ne} \PY{o}{=} \PY{n}{np}\PY{o}{.}\PY{n}{arange}\PY{p}{(}\PY{n}{ne\PYZus{}gray}\PY{o}{.}\PY{n}{shape}\PY{p}{[}\PY{l+m+mi}{1}\PY{p}{]}\PY{p}{)}
\PY{n}{trace\PYZus{}center\PYZus{}ne} \PY{o}{=} \PY{n}{trace\PYZus{}model}\PY{p}{(}\PY{n}{xaxis\PYZus{}ne}\PY{p}{)}
\end{Verbatim}
\end{tcolorbox}

    \begin{tcolorbox}[breakable, size=fbox, boxrule=1pt, pad at break*=1mm,colback=cellbackground, colframe=cellborder]
\prompt{In}{incolor}{92}{\boxspacing}
\begin{Verbatim}[commandchars=\\\{\}]
\PY{n}{xaxis\PYZus{}kr} \PY{o}{=} \PY{n}{np}\PY{o}{.}\PY{n}{arange}\PY{p}{(}\PY{n}{kr\PYZus{}gray}\PY{o}{.}\PY{n}{shape}\PY{p}{[}\PY{l+m+mi}{1}\PY{p}{]}\PY{p}{)}
\PY{n}{trace\PYZus{}center\PYZus{}kr} \PY{o}{=} \PY{n}{trace\PYZus{}model}\PY{p}{(}\PY{n}{xaxis\PYZus{}kr}\PY{p}{)}
\end{Verbatim}
\end{tcolorbox}

    \begin{tcolorbox}[breakable, size=fbox, boxrule=1pt, pad at break*=1mm,colback=cellbackground, colframe=cellborder]
\prompt{In}{incolor}{93}{\boxspacing}
\begin{Verbatim}[commandchars=\\\{\}]
\PY{n}{xaxis\PYZus{}he} \PY{o}{=} \PY{n}{np}\PY{o}{.}\PY{n}{arange}\PY{p}{(}\PY{n}{he\PYZus{}gray}\PY{o}{.}\PY{n}{shape}\PY{p}{[}\PY{l+m+mi}{1}\PY{p}{]}\PY{p}{)}
\PY{n}{trace\PYZus{}center\PYZus{}he} \PY{o}{=} \PY{n}{trace\PYZus{}model}\PY{p}{(}\PY{n}{xaxis\PYZus{}he}\PY{p}{)}
\end{Verbatim}
\end{tcolorbox}

    Luego, usamos la traza para extraer espectros de cada uno de los tres
espectros de las lámparas de calibración.

Primero verificamos que las trazas se vean aceptables, ajustamos los
límites de los ejes x y y para poder observar el espectro y las trazas
de la mejor manera posible

    \begin{tcolorbox}[breakable, size=fbox, boxrule=1pt, pad at break*=1mm,colback=cellbackground, colframe=cellborder]
\prompt{In}{incolor}{98}{\boxspacing}
\begin{Verbatim}[commandchars=\\\{\}]
\PY{n}{plt}\PY{o}{.}\PY{n}{figure}\PY{p}{(}\PY{n}{figsize}\PY{o}{=}\PY{p}{(}\PY{l+m+mi}{16}\PY{p}{,}\PY{l+m+mi}{8}\PY{p}{)}\PY{p}{)}
\PY{n}{ax1} \PY{o}{=} \PY{n}{plt}\PY{o}{.}\PY{n}{subplot}\PY{p}{(}\PY{l+m+mi}{4}\PY{p}{,}\PY{l+m+mi}{1}\PY{p}{,}\PY{l+m+mi}{1}\PY{p}{)}
\PY{n}{ax1}\PY{o}{.}\PY{n}{imshow}\PY{p}{(}\PY{n}{hg\PYZus{}gray}\PY{p}{[}\PY{l+m+mi}{250}\PY{p}{:}\PY{l+m+mi}{300}\PY{p}{,}\PY{p}{:}\PY{p}{]}\PY{p}{,} \PY{n}{extent}\PY{o}{=}\PY{p}{[}\PY{l+m+mi}{0}\PY{p}{,}\PY{l+m+mi}{1000}\PY{p}{,}\PY{l+m+mi}{250}\PY{p}{,}\PY{l+m+mi}{350}\PY{p}{]}\PY{p}{)}
\PY{n}{ax1}\PY{o}{.}\PY{n}{plot}\PY{p}{(}\PY{n}{xaxis\PYZus{}he}\PY{p}{,} \PY{n}{trace\PYZus{}center\PYZus{}he}\PY{p}{)}
\PY{n}{ax1}\PY{o}{.}\PY{n}{set\PYZus{}aspect}\PY{p}{(}\PY{l+m+mi}{3}\PY{p}{)}
\PY{n}{ax2} \PY{o}{=} \PY{n}{plt}\PY{o}{.}\PY{n}{subplot}\PY{p}{(}\PY{l+m+mi}{4}\PY{p}{,}\PY{l+m+mi}{1}\PY{p}{,}\PY{l+m+mi}{2}\PY{p}{)}
\PY{n}{ax2}\PY{o}{.}\PY{n}{imshow}\PY{p}{(}\PY{n}{kr\PYZus{}gray}\PY{p}{[}\PY{l+m+mi}{250}\PY{p}{:}\PY{l+m+mi}{300}\PY{p}{,}\PY{p}{:}\PY{p}{]}\PY{p}{,} \PY{n}{extent}\PY{o}{=}\PY{p}{[}\PY{l+m+mi}{0}\PY{p}{,}\PY{l+m+mi}{1000}\PY{p}{,}\PY{l+m+mi}{250}\PY{p}{,}\PY{l+m+mi}{350}\PY{p}{]}\PY{p}{)}
\PY{n}{ax2}\PY{o}{.}\PY{n}{plot}\PY{p}{(}\PY{n}{xaxis\PYZus{}he}\PY{p}{,} \PY{n}{trace\PYZus{}center\PYZus{}he}\PY{p}{)}
\PY{n}{ax2}\PY{o}{.}\PY{n}{set\PYZus{}aspect}\PY{p}{(}\PY{l+m+mi}{3}\PY{p}{)}
\PY{n}{ax3} \PY{o}{=} \PY{n}{plt}\PY{o}{.}\PY{n}{subplot}\PY{p}{(}\PY{l+m+mi}{4}\PY{p}{,}\PY{l+m+mi}{1}\PY{p}{,}\PY{l+m+mi}{3}\PY{p}{)}
\PY{n}{ax3}\PY{o}{.}\PY{n}{imshow}\PY{p}{(}\PY{n}{ne\PYZus{}gray}\PY{p}{[}\PY{l+m+mi}{250}\PY{p}{:}\PY{l+m+mi}{300}\PY{p}{,}\PY{p}{:}\PY{p}{]}\PY{p}{,} \PY{n}{extent}\PY{o}{=}\PY{p}{[}\PY{l+m+mi}{0}\PY{p}{,}\PY{l+m+mi}{1000}\PY{p}{,}\PY{l+m+mi}{250}\PY{p}{,}\PY{l+m+mi}{350}\PY{p}{]}\PY{p}{)}
\PY{n}{ax3}\PY{o}{.}\PY{n}{plot}\PY{p}{(}\PY{n}{xaxis\PYZus{}he}\PY{p}{,} \PY{n}{trace\PYZus{}center\PYZus{}he}\PY{p}{)}
\PY{n}{ax3}\PY{o}{.}\PY{n}{set\PYZus{}aspect}\PY{p}{(}\PY{l+m+mi}{3}\PY{p}{)}
\PY{n}{ax4} \PY{o}{=} \PY{n}{plt}\PY{o}{.}\PY{n}{subplot}\PY{p}{(}\PY{l+m+mi}{4}\PY{p}{,}\PY{l+m+mi}{1}\PY{p}{,}\PY{l+m+mi}{4}\PY{p}{)}
\PY{n}{ax4}\PY{o}{.}\PY{n}{imshow}\PY{p}{(}\PY{n}{he\PYZus{}gray}\PY{p}{[}\PY{l+m+mi}{250}\PY{p}{:}\PY{l+m+mi}{300}\PY{p}{,}\PY{p}{:}\PY{p}{]}\PY{p}{,} \PY{n}{extent}\PY{o}{=}\PY{p}{[}\PY{l+m+mi}{0}\PY{p}{,}\PY{l+m+mi}{1000}\PY{p}{,}\PY{l+m+mi}{250}\PY{p}{,}\PY{l+m+mi}{350}\PY{p}{]}\PY{p}{)}
\PY{n}{ax4}\PY{o}{.}\PY{n}{plot}\PY{p}{(}\PY{n}{xaxis\PYZus{}he}\PY{p}{,} \PY{n}{trace\PYZus{}center\PYZus{}he}\PY{p}{)}
\PY{n}{ax4}\PY{o}{.}\PY{n}{set\PYZus{}aspect}\PY{p}{(}\PY{l+m+mi}{3}\PY{p}{)}
\end{Verbatim}
\end{tcolorbox}

    \begin{center}
    \adjustimage{max size={0.9\linewidth}{0.9\paperheight}}{output_30_0.png}
    \end{center}
    { \hspace*{\fill} \\}
    
    Mostramos individualmente cada espectro con su respectiva traza, para
cada uno de los elementos empleados:

    \begin{tcolorbox}[breakable, size=fbox, boxrule=1pt, pad at break*=1mm,colback=cellbackground, colframe=cellborder]
\prompt{In}{incolor}{99}{\boxspacing}
\begin{Verbatim}[commandchars=\\\{\}]
\PY{n}{plt}\PY{o}{.}\PY{n}{figure}\PY{p}{(}\PY{n}{figsize}\PY{o}{=}\PY{p}{(}\PY{l+m+mi}{16}\PY{p}{,}\PY{l+m+mi}{8}\PY{p}{)}\PY{p}{)}
\PY{n}{ax1} \PY{o}{=} \PY{n}{plt}\PY{o}{.}\PY{n}{subplot}\PY{p}{(}\PY{l+m+mi}{1}\PY{p}{,}\PY{l+m+mi}{1}\PY{p}{,}\PY{l+m+mi}{1}\PY{p}{)}
\PY{n}{ax1}\PY{o}{.}\PY{n}{imshow}\PY{p}{(}\PY{n}{hg\PYZus{}gray}\PY{p}{[}\PY{l+m+mi}{50}\PY{p}{:}\PY{l+m+mi}{150}\PY{p}{,}\PY{p}{:}\PY{p}{]}\PY{p}{,} \PY{n}{extent}\PY{o}{=}\PY{p}{[}\PY{l+m+mi}{0}\PY{p}{,} \PY{l+m+mi}{600}\PY{p}{,}\PY{l+m+mi}{250}\PY{p}{,}\PY{l+m+mi}{350}\PY{p}{]}\PY{p}{)}
\PY{n}{ax1}\PY{o}{.}\PY{n}{plot}\PY{p}{(}\PY{n}{xaxis\PYZus{}hg}\PY{p}{,} \PY{n}{trace\PYZus{}center\PYZus{}hg}\PY{p}{)}
\end{Verbatim}
\end{tcolorbox}

            \begin{tcolorbox}[breakable, size=fbox, boxrule=.5pt, pad at break*=1mm, opacityfill=0]
\prompt{Out}{outcolor}{99}{\boxspacing}
\begin{Verbatim}[commandchars=\\\{\}]
[<matplotlib.lines.Line2D at 0x7f2f3421ffd0>]
\end{Verbatim}
\end{tcolorbox}
        
    \begin{center}
    \adjustimage{max size={0.9\linewidth}{0.9\paperheight}}{output_32_1.png}
    \end{center}
    { \hspace*{\fill} \\}
    
    \begin{tcolorbox}[breakable, size=fbox, boxrule=1pt, pad at break*=1mm,colback=cellbackground, colframe=cellborder]
\prompt{In}{incolor}{100}{\boxspacing}
\begin{Verbatim}[commandchars=\\\{\}]
\PY{n}{plt}\PY{o}{.}\PY{n}{figure}\PY{p}{(}\PY{n}{figsize}\PY{o}{=}\PY{p}{(}\PY{l+m+mi}{16}\PY{p}{,}\PY{l+m+mi}{8}\PY{p}{)}\PY{p}{)}
\PY{n}{ax1} \PY{o}{=} \PY{n}{plt}\PY{o}{.}\PY{n}{subplot}\PY{p}{(}\PY{l+m+mi}{1}\PY{p}{,}\PY{l+m+mi}{1}\PY{p}{,}\PY{l+m+mi}{1}\PY{p}{)}
\PY{n}{ax1}\PY{o}{.}\PY{n}{imshow}\PY{p}{(}\PY{n}{ne\PYZus{}gray}\PY{p}{[}\PY{l+m+mi}{250}\PY{p}{:}\PY{l+m+mi}{300}\PY{p}{,}\PY{p}{:}\PY{p}{]}\PY{p}{,} \PY{n}{extent}\PY{o}{=}\PY{p}{[}\PY{l+m+mi}{0}\PY{p}{,} \PY{l+m+mi}{950}\PY{p}{,}\PY{l+m+mi}{250}\PY{p}{,}\PY{l+m+mi}{350}\PY{p}{]}\PY{p}{)}
\PY{n}{ax1}\PY{o}{.}\PY{n}{plot}\PY{p}{(}\PY{n}{xaxis\PYZus{}ne}\PY{p}{,} \PY{n}{trace\PYZus{}center\PYZus{}ne}\PY{p}{)}
\end{Verbatim}
\end{tcolorbox}

            \begin{tcolorbox}[breakable, size=fbox, boxrule=.5pt, pad at break*=1mm, opacityfill=0]
\prompt{Out}{outcolor}{100}{\boxspacing}
\begin{Verbatim}[commandchars=\\\{\}]
[<matplotlib.lines.Line2D at 0x7f2f342e0b90>]
\end{Verbatim}
\end{tcolorbox}
        
    \begin{center}
    \adjustimage{max size={0.9\linewidth}{0.9\paperheight}}{output_33_1.png}
    \end{center}
    { \hspace*{\fill} \\}
    
    \begin{tcolorbox}[breakable, size=fbox, boxrule=1pt, pad at break*=1mm,colback=cellbackground, colframe=cellborder]
\prompt{In}{incolor}{101}{\boxspacing}
\begin{Verbatim}[commandchars=\\\{\}]
\PY{n}{plt}\PY{o}{.}\PY{n}{figure}\PY{p}{(}\PY{n}{figsize}\PY{o}{=}\PY{p}{(}\PY{l+m+mi}{16}\PY{p}{,}\PY{l+m+mi}{8}\PY{p}{)}\PY{p}{)}
\PY{n}{ax1} \PY{o}{=} \PY{n}{plt}\PY{o}{.}\PY{n}{subplot}\PY{p}{(}\PY{l+m+mi}{1}\PY{p}{,}\PY{l+m+mi}{1}\PY{p}{,}\PY{l+m+mi}{1}\PY{p}{)}
\PY{n}{ax1}\PY{o}{.}\PY{n}{imshow}\PY{p}{(}\PY{n}{kr\PYZus{}gray}\PY{p}{[}\PY{l+m+mi}{250}\PY{p}{:}\PY{l+m+mi}{300}\PY{p}{,}\PY{p}{:}\PY{p}{]}\PY{p}{,} \PY{n}{extent}\PY{o}{=}\PY{p}{[}\PY{l+m+mi}{0}\PY{p}{,} \PY{l+m+mi}{600}\PY{p}{,}\PY{l+m+mi}{250}\PY{p}{,}\PY{l+m+mi}{350}\PY{p}{]}\PY{p}{)}
\PY{n}{ax1}\PY{o}{.}\PY{n}{plot}\PY{p}{(}\PY{n}{xaxis\PYZus{}kr}\PY{p}{,} \PY{n}{trace\PYZus{}center\PYZus{}kr}\PY{p}{)}
\end{Verbatim}
\end{tcolorbox}

            \begin{tcolorbox}[breakable, size=fbox, boxrule=.5pt, pad at break*=1mm, opacityfill=0]
\prompt{Out}{outcolor}{101}{\boxspacing}
\begin{Verbatim}[commandchars=\\\{\}]
[<matplotlib.lines.Line2D at 0x7f2f34172410>]
\end{Verbatim}
\end{tcolorbox}
        
    \begin{center}
    \adjustimage{max size={0.9\linewidth}{0.9\paperheight}}{output_34_1.png}
    \end{center}
    { \hspace*{\fill} \\}
    
    \begin{tcolorbox}[breakable, size=fbox, boxrule=1pt, pad at break*=1mm,colback=cellbackground, colframe=cellborder]
\prompt{In}{incolor}{102}{\boxspacing}
\begin{Verbatim}[commandchars=\\\{\}]
\PY{n}{plt}\PY{o}{.}\PY{n}{figure}\PY{p}{(}\PY{n}{figsize}\PY{o}{=}\PY{p}{(}\PY{l+m+mi}{16}\PY{p}{,}\PY{l+m+mi}{8}\PY{p}{)}\PY{p}{)}
\PY{n}{ax1} \PY{o}{=} \PY{n}{plt}\PY{o}{.}\PY{n}{subplot}\PY{p}{(}\PY{l+m+mi}{1}\PY{p}{,}\PY{l+m+mi}{1}\PY{p}{,}\PY{l+m+mi}{1}\PY{p}{)}
\PY{n}{ax1}\PY{o}{.}\PY{n}{imshow}\PY{p}{(}\PY{n}{he\PYZus{}gray}\PY{p}{[}\PY{l+m+mi}{250}\PY{p}{:}\PY{l+m+mi}{300}\PY{p}{,}\PY{p}{:}\PY{p}{]}\PY{p}{,} \PY{n}{extent}\PY{o}{=}\PY{p}{[}\PY{l+m+mi}{0}\PY{p}{,} \PY{l+m+mi}{1000}\PY{p}{,}\PY{l+m+mi}{250}\PY{p}{,}\PY{l+m+mi}{350}\PY{p}{]}\PY{p}{)}
\PY{n}{ax1}\PY{o}{.}\PY{n}{plot}\PY{p}{(}\PY{n}{xaxis\PYZus{}he}\PY{p}{,} \PY{n}{trace\PYZus{}center\PYZus{}he}\PY{p}{)}
\end{Verbatim}
\end{tcolorbox}

            \begin{tcolorbox}[breakable, size=fbox, boxrule=.5pt, pad at break*=1mm, opacityfill=0]
\prompt{Out}{outcolor}{102}{\boxspacing}
\begin{Verbatim}[commandchars=\\\{\}]
[<matplotlib.lines.Line2D at 0x7f2f34198550>]
\end{Verbatim}
\end{tcolorbox}
        
    \begin{center}
    \adjustimage{max size={0.9\linewidth}{0.9\paperheight}}{output_35_1.png}
    \end{center}
    { \hspace*{\fill} \\}
    
    En las siguientes tres celdas, extraemos el espectro promedio ponderado
por el perfil de la traza de cada uno de los diferentes elementos que se
tienen

\texttt{npixels\_to\_cut} establece una región recortada alrededor de la
traza, similar a lo que se muestra en las figuras anteriores. El
\texttt{model\_trace\_profile}, derivado anteriormente, es el perfil de
transmisión que medimos usando la traza estelar.

    \begin{tcolorbox}[breakable, size=fbox, boxrule=1pt, pad at break*=1mm,colback=cellbackground, colframe=cellborder]
\prompt{In}{incolor}{103}{\boxspacing}
\begin{Verbatim}[commandchars=\\\{\}]
\PY{n}{hg\PYZus{}spectrum} \PY{o}{=} \PY{n}{np}\PY{o}{.}\PY{n}{array}\PY{p}{(}\PY{p}{[}\PY{n}{np}\PY{o}{.}\PY{n}{average}\PY{p}{(}\PY{n}{hg\PYZus{}gray}\PY{p}{[}\PY{n+nb}{int}\PY{p}{(}\PY{n}{yval}\PY{p}{)}\PY{o}{\PYZhy{}}\PY{n}{npixels\PYZus{}to\PYZus{}cut}\PY{p}{:}\PY{n+nb}{int}\PY{p}{(}\PY{n}{yval}\PY{p}{)}\PY{o}{+}\PY{n}{npixels\PYZus{}to\PYZus{}cut}\PY{p}{,} \PY{n}{ii}\PY{p}{]}\PY{p}{,}
                                 \PY{n}{weights}\PY{o}{=}\PY{n}{model\PYZus{}trace\PYZus{}profile}\PY{p}{)}
                      \PY{k}{for} \PY{n}{yval}\PY{p}{,} \PY{n}{ii} \PY{o+ow}{in} \PY{n+nb}{zip}\PY{p}{(}\PY{n}{trace\PYZus{}center\PYZus{}hg}\PY{p}{,} \PY{n}{xaxis\PYZus{}hg}\PY{p}{)}\PY{p}{]}\PY{p}{)}
\end{Verbatim}
\end{tcolorbox}

    \begin{tcolorbox}[breakable, size=fbox, boxrule=1pt, pad at break*=1mm,colback=cellbackground, colframe=cellborder]
\prompt{In}{incolor}{104}{\boxspacing}
\begin{Verbatim}[commandchars=\\\{\}]
\PY{n}{ne\PYZus{}spectrum} \PY{o}{=} \PY{n}{np}\PY{o}{.}\PY{n}{array}\PY{p}{(}\PY{p}{[}\PY{n}{np}\PY{o}{.}\PY{n}{average}\PY{p}{(}\PY{n}{ne\PYZus{}gray}\PY{p}{[}\PY{n+nb}{int}\PY{p}{(}\PY{n}{yval}\PY{p}{)}\PY{o}{\PYZhy{}}\PY{n}{npixels\PYZus{}to\PYZus{}cut}\PY{p}{:}\PY{n+nb}{int}\PY{p}{(}\PY{n}{yval}\PY{p}{)}\PY{o}{+}\PY{n}{npixels\PYZus{}to\PYZus{}cut}\PY{p}{,} \PY{n}{ii}\PY{p}{]}\PY{p}{,}
                                \PY{n}{weights}\PY{o}{=}\PY{n}{model\PYZus{}trace\PYZus{}profile}\PY{p}{)}
                     \PY{k}{for} \PY{n}{yval}\PY{p}{,} \PY{n}{ii} \PY{o+ow}{in} \PY{n+nb}{zip}\PY{p}{(}\PY{n}{trace\PYZus{}center\PYZus{}ne}\PY{p}{,} \PY{n}{xaxis\PYZus{}ne}\PY{p}{)}\PY{p}{]}\PY{p}{)}
\end{Verbatim}
\end{tcolorbox}

    \begin{tcolorbox}[breakable, size=fbox, boxrule=1pt, pad at break*=1mm,colback=cellbackground, colframe=cellborder]
\prompt{In}{incolor}{105}{\boxspacing}
\begin{Verbatim}[commandchars=\\\{\}]
\PY{n}{kr\PYZus{}spectrum} \PY{o}{=} \PY{n}{np}\PY{o}{.}\PY{n}{array}\PY{p}{(}\PY{p}{[}\PY{n}{np}\PY{o}{.}\PY{n}{average}\PY{p}{(}\PY{n}{kr\PYZus{}gray}\PY{p}{[}\PY{n+nb}{int}\PY{p}{(}\PY{n}{yval}\PY{p}{)}\PY{o}{\PYZhy{}}\PY{n}{npixels\PYZus{}to\PYZus{}cut}\PY{p}{:}\PY{n+nb}{int}\PY{p}{(}\PY{n}{yval}\PY{p}{)}\PY{o}{+}\PY{n}{npixels\PYZus{}to\PYZus{}cut}\PY{p}{,} \PY{n}{ii}\PY{p}{]}\PY{p}{,}
                                \PY{n}{weights}\PY{o}{=}\PY{n}{model\PYZus{}trace\PYZus{}profile}\PY{p}{)}
                     \PY{k}{for} \PY{n}{yval}\PY{p}{,} \PY{n}{ii} \PY{o+ow}{in} \PY{n+nb}{zip}\PY{p}{(}\PY{n}{trace\PYZus{}center\PYZus{}kr}\PY{p}{,} \PY{n}{xaxis\PYZus{}kr}\PY{p}{)}\PY{p}{]}\PY{p}{)}
\end{Verbatim}
\end{tcolorbox}

    \begin{tcolorbox}[breakable, size=fbox, boxrule=1pt, pad at break*=1mm,colback=cellbackground, colframe=cellborder]
\prompt{In}{incolor}{106}{\boxspacing}
\begin{Verbatim}[commandchars=\\\{\}]
\PY{n}{he\PYZus{}spectrum} \PY{o}{=} \PY{n}{np}\PY{o}{.}\PY{n}{array}\PY{p}{(}\PY{p}{[}\PY{n}{np}\PY{o}{.}\PY{n}{average}\PY{p}{(}\PY{n}{he\PYZus{}gray}\PY{p}{[}\PY{n+nb}{int}\PY{p}{(}\PY{n}{yval}\PY{p}{)}\PY{o}{\PYZhy{}}\PY{n}{npixels\PYZus{}to\PYZus{}cut}\PY{p}{:}\PY{n+nb}{int}\PY{p}{(}\PY{n}{yval}\PY{p}{)}\PY{o}{+}\PY{n}{npixels\PYZus{}to\PYZus{}cut}\PY{p}{,} \PY{n}{ii}\PY{p}{]}\PY{p}{,}
                                   \PY{n}{weights}\PY{o}{=}\PY{n}{model\PYZus{}trace\PYZus{}profile}\PY{p}{)}
                    \PY{k}{for} \PY{n}{yval}\PY{p}{,} \PY{n}{ii} \PY{o+ow}{in} \PY{n+nb}{zip}\PY{p}{(}\PY{n}{trace\PYZus{}center\PYZus{}he}\PY{p}{,} \PY{n}{xaxis\PYZus{}he}\PY{p}{)}\PY{p}{]}\PY{p}{)}
\end{Verbatim}
\end{tcolorbox}

    Se observan los espectros de los elementos

    \begin{tcolorbox}[breakable, size=fbox, boxrule=1pt, pad at break*=1mm,colback=cellbackground, colframe=cellborder]
\prompt{In}{incolor}{107}{\boxspacing}
\begin{Verbatim}[commandchars=\\\{\}]
\PY{n}{plt}\PY{o}{.}\PY{n}{plot}\PY{p}{(}\PY{n}{xaxis\PYZus{}hg}\PY{p}{,} \PY{n}{hg\PYZus{}spectrum}\PY{p}{,} \PY{n}{label}\PY{o}{=}\PY{l+s+s1}{\PYZsq{}}\PY{l+s+s1}{Mercury}\PY{l+s+s1}{\PYZsq{}}\PY{p}{)}
\PY{n}{plt}\PY{o}{.}\PY{n}{plot}\PY{p}{(}\PY{n}{xaxis\PYZus{}ne}\PY{p}{,} \PY{n}{ne\PYZus{}spectrum}\PY{p}{,} \PY{n}{label}\PY{o}{=}\PY{l+s+s1}{\PYZsq{}}\PY{l+s+s1}{Neon}\PY{l+s+s1}{\PYZsq{}}\PY{p}{)}
\PY{n}{plt}\PY{o}{.}\PY{n}{plot}\PY{p}{(}\PY{n}{xaxis\PYZus{}kr}\PY{p}{,} \PY{n}{kr\PYZus{}spectrum}\PY{p}{,} \PY{n}{label}\PY{o}{=}\PY{l+s+s1}{\PYZsq{}}\PY{l+s+s1}{Krypton}\PY{l+s+s1}{\PYZsq{}}\PY{p}{)}
\PY{n}{plt}\PY{o}{.}\PY{n}{plot}\PY{p}{(}\PY{n}{xaxis\PYZus{}he}\PY{p}{,} \PY{n}{he\PYZus{}spectrum}\PY{p}{,} \PY{n}{label}\PY{o}{=}\PY{l+s+s1}{\PYZsq{}}\PY{l+s+s1}{Helium}\PY{l+s+s1}{\PYZsq{}}\PY{p}{)}
\PY{n}{plt}\PY{o}{.}\PY{n}{legend}\PY{p}{(}\PY{n}{loc}\PY{o}{=}\PY{l+s+s1}{\PYZsq{}}\PY{l+s+s1}{best}\PY{l+s+s1}{\PYZsq{}}\PY{p}{)}\PY{p}{;}
\end{Verbatim}
\end{tcolorbox}

    \begin{center}
    \adjustimage{max size={0.9\linewidth}{0.9\paperheight}}{output_42_0.png}
    \end{center}
    { \hspace*{\fill} \\}
    
    \begin{tcolorbox}[breakable, size=fbox, boxrule=1pt, pad at break*=1mm,colback=cellbackground, colframe=cellborder]
\prompt{In}{incolor}{108}{\boxspacing}
\begin{Verbatim}[commandchars=\\\{\}]
\PY{k+kn}{from} \PY{n+nn}{IPython}\PY{n+nn}{.}\PY{n+nn}{display} \PY{k+kn}{import} \PY{n}{Image}\PY{p}{,} \PY{n}{display}
\PY{k+kn}{import} \PY{n+nn}{requests}
\end{Verbatim}
\end{tcolorbox}

    Image(``https://upload.wikimedia.org/wikipedia/commons/2/29/Mercury\_Spectra.jpg'')

    Image(``https://upload.wikimedia.org/wikipedia/commons/9/99/Neon\_spectra.jpg'')

    Image(``https://upload.wikimedia.org/wikipedia/commons/a/a6/Krypton\_Spectrum.jpg'')

    \begin{tcolorbox}[breakable, size=fbox, boxrule=1pt, pad at break*=1mm,colback=cellbackground, colframe=cellborder]
\prompt{In}{incolor}{16}{\boxspacing}
\begin{Verbatim}[commandchars=\\\{\}]
\PY{n}{url} \PY{o}{=} \PY{l+s+s2}{\PYZdq{}}\PY{l+s+s2}{https://upload.wikimedia.org/wikipedia/commons/2/29/Mercury\PYZus{}Spectra.jpg}\PY{l+s+s2}{\PYZdq{}}
\PY{n}{headers} \PY{o}{=} \PY{p}{\PYZob{}}\PY{l+s+s2}{\PYZdq{}}\PY{l+s+s2}{User\PYZhy{}Agent}\PY{l+s+s2}{\PYZdq{}}\PY{p}{:} \PY{l+s+s2}{\PYZdq{}}\PY{l+s+s2}{Mozilla/5.0}\PY{l+s+s2}{\PYZdq{}}\PY{p}{\PYZcb{}}

\PY{n}{r} \PY{o}{=} \PY{n}{requests}\PY{o}{.}\PY{n}{get}\PY{p}{(}\PY{n}{url}\PY{p}{,} \PY{n}{headers}\PY{o}{=}\PY{n}{headers}\PY{p}{,} \PY{n}{timeout}\PY{o}{=}\PY{l+m+mi}{15}\PY{p}{)}
\PY{n}{r}\PY{o}{.}\PY{n}{raise\PYZus{}for\PYZus{}status}\PY{p}{(}\PY{p}{)}              \PY{c+c1}{\PYZsh{} lanza error si no es 200}
\PY{n}{display}\PY{p}{(}\PY{n}{Image}\PY{p}{(}\PY{n}{data}\PY{o}{=}\PY{n}{r}\PY{o}{.}\PY{n}{content}\PY{p}{)}\PY{p}{)}    \PY{c+c1}{\PYZsh{} muestra el binario directamente}
\end{Verbatim}
\end{tcolorbox}

    \begin{center}
    \adjustimage{max size={0.9\linewidth}{0.9\paperheight}}{output_47_0.jpg}
    \end{center}
    { \hspace*{\fill} \\}
    
    \begin{tcolorbox}[breakable, size=fbox, boxrule=1pt, pad at break*=1mm,colback=cellbackground, colframe=cellborder]
\prompt{In}{incolor}{17}{\boxspacing}
\begin{Verbatim}[commandchars=\\\{\}]
\PY{n}{url} \PY{o}{=} \PY{l+s+s2}{\PYZdq{}}\PY{l+s+s2}{https://upload.wikimedia.org/wikipedia/commons/9/99/Neon\PYZus{}spectra.jpg}\PY{l+s+s2}{\PYZdq{}}
\PY{n}{headers} \PY{o}{=} \PY{p}{\PYZob{}}\PY{l+s+s2}{\PYZdq{}}\PY{l+s+s2}{User\PYZhy{}Agent}\PY{l+s+s2}{\PYZdq{}}\PY{p}{:} \PY{l+s+s2}{\PYZdq{}}\PY{l+s+s2}{Mozilla/5.0}\PY{l+s+s2}{\PYZdq{}}\PY{p}{\PYZcb{}}

\PY{n}{r} \PY{o}{=} \PY{n}{requests}\PY{o}{.}\PY{n}{get}\PY{p}{(}\PY{n}{url}\PY{p}{,} \PY{n}{headers}\PY{o}{=}\PY{n}{headers}\PY{p}{,} \PY{n}{timeout}\PY{o}{=}\PY{l+m+mi}{15}\PY{p}{)}
\PY{n}{r}\PY{o}{.}\PY{n}{raise\PYZus{}for\PYZus{}status}\PY{p}{(}\PY{p}{)}              \PY{c+c1}{\PYZsh{} lanza error si no es 200}
\PY{n}{display}\PY{p}{(}\PY{n}{Image}\PY{p}{(}\PY{n}{data}\PY{o}{=}\PY{n}{r}\PY{o}{.}\PY{n}{content}\PY{p}{)}\PY{p}{)}    \PY{c+c1}{\PYZsh{} muestra el binario directamente}
\end{Verbatim}
\end{tcolorbox}

    \begin{center}
    \adjustimage{max size={0.9\linewidth}{0.9\paperheight}}{output_48_0.jpg}
    \end{center}
    { \hspace*{\fill} \\}
    
    \begin{tcolorbox}[breakable, size=fbox, boxrule=1pt, pad at break*=1mm,colback=cellbackground, colframe=cellborder]
\prompt{In}{incolor}{18}{\boxspacing}
\begin{Verbatim}[commandchars=\\\{\}]
\PY{n}{url} \PY{o}{=} \PY{l+s+s2}{\PYZdq{}}\PY{l+s+s2}{https://upload.wikimedia.org/wikipedia/commons/a/a6/Krypton\PYZus{}Spectrum.jpg}\PY{l+s+s2}{\PYZdq{}}
\PY{n}{headers} \PY{o}{=} \PY{p}{\PYZob{}}\PY{l+s+s2}{\PYZdq{}}\PY{l+s+s2}{User\PYZhy{}Agent}\PY{l+s+s2}{\PYZdq{}}\PY{p}{:} \PY{l+s+s2}{\PYZdq{}}\PY{l+s+s2}{Mozilla/5.0}\PY{l+s+s2}{\PYZdq{}}\PY{p}{\PYZcb{}}

\PY{n}{r} \PY{o}{=} \PY{n}{requests}\PY{o}{.}\PY{n}{get}\PY{p}{(}\PY{n}{url}\PY{p}{,} \PY{n}{headers}\PY{o}{=}\PY{n}{headers}\PY{p}{,} \PY{n}{timeout}\PY{o}{=}\PY{l+m+mi}{15}\PY{p}{)}
\PY{n}{r}\PY{o}{.}\PY{n}{raise\PYZus{}for\PYZus{}status}\PY{p}{(}\PY{p}{)}              \PY{c+c1}{\PYZsh{} lanza error si no es 200}
\PY{n}{display}\PY{p}{(}\PY{n}{Image}\PY{p}{(}\PY{n}{data}\PY{o}{=}\PY{n}{r}\PY{o}{.}\PY{n}{content}\PY{p}{)}\PY{p}{)}    \PY{c+c1}{\PYZsh{} muestra el binario directamente}
\end{Verbatim}
\end{tcolorbox}

    \begin{center}
    \adjustimage{max size={0.9\linewidth}{0.9\paperheight}}{output_49_0.jpg}
    \end{center}
    { \hspace*{\fill} \\}
    
    Image(``https://upload.wikimedia.org/wikipedia/commons/0/0d/HG-Spektrum\_crop.jpg'')

    \begin{tcolorbox}[breakable, size=fbox, boxrule=1pt, pad at break*=1mm,colback=cellbackground, colframe=cellborder]
\prompt{In}{incolor}{19}{\boxspacing}
\begin{Verbatim}[commandchars=\\\{\}]
\PY{n}{url} \PY{o}{=} \PY{l+s+s2}{\PYZdq{}}\PY{l+s+s2}{https://upload.wikimedia.org/wikipedia/commons/0/0d/HG\PYZhy{}Spektrum\PYZus{}crop.jpg}\PY{l+s+s2}{\PYZdq{}}
\PY{n}{headers} \PY{o}{=} \PY{p}{\PYZob{}}\PY{l+s+s2}{\PYZdq{}}\PY{l+s+s2}{User\PYZhy{}Agent}\PY{l+s+s2}{\PYZdq{}}\PY{p}{:} \PY{l+s+s2}{\PYZdq{}}\PY{l+s+s2}{Mozilla/5.0}\PY{l+s+s2}{\PYZdq{}}\PY{p}{\PYZcb{}}

\PY{n}{r} \PY{o}{=} \PY{n}{requests}\PY{o}{.}\PY{n}{get}\PY{p}{(}\PY{n}{url}\PY{p}{,} \PY{n}{headers}\PY{o}{=}\PY{n}{headers}\PY{p}{,} \PY{n}{timeout}\PY{o}{=}\PY{l+m+mi}{15}\PY{p}{)}
\PY{n}{r}\PY{o}{.}\PY{n}{raise\PYZus{}for\PYZus{}status}\PY{p}{(}\PY{p}{)}              \PY{c+c1}{\PYZsh{} lanza error si no es 200}
\PY{n}{display}\PY{p}{(}\PY{n}{Image}\PY{p}{(}\PY{n}{data}\PY{o}{=}\PY{n}{r}\PY{o}{.}\PY{n}{content}\PY{p}{)}\PY{p}{)}    \PY{c+c1}{\PYZsh{} muestra el binario directamente}
\end{Verbatim}
\end{tcolorbox}

    \begin{center}
    \adjustimage{max size={0.9\linewidth}{0.9\paperheight}}{output_51_0.jpg}
    \end{center}
    { \hspace*{\fill} \\}
    
    The optical lines of Mercury are, from Wikipedia, in Angstroms:

\begin{itemize}
\tightlist
\item
  4047 violet
\item
  4358 blue
\item
  5461 green
\item
  5782 yellow-orange
\item
  6500 red
\end{itemize}

    Se imprimen individualmente los espectros del Mercurio, el Neón, el
Kriptón y el Helio respectivamente, se escogen los picos que van a ser
comparados con los valores teóricos sacados del NIST o de alguna base de
datos que nos permita saber las líneas espectrales características de
cada element químico. Tenemos que tener en cuenta los pixeles en los que
están ubicados estos picos que vamos a comparar.

    \textbf{\emph{Mercurio}}

    \begin{tcolorbox}[breakable, size=fbox, boxrule=1pt, pad at break*=1mm,colback=cellbackground, colframe=cellborder]
\prompt{In}{incolor}{110}{\boxspacing}
\begin{Verbatim}[commandchars=\\\{\}]
\PY{n}{plt}\PY{o}{.}\PY{n}{plot}\PY{p}{(}\PY{n}{xaxis\PYZus{}hg}\PY{p}{,} \PY{n}{hg\PYZus{}spectrum}\PY{p}{)}\PY{p}{;}
\end{Verbatim}
\end{tcolorbox}

    \begin{center}
    \adjustimage{max size={0.9\linewidth}{0.9\paperheight}}{output_55_0.png}
    \end{center}
    { \hspace*{\fill} \\}
    
    \textbf{\emph{Neón}}

    \begin{tcolorbox}[breakable, size=fbox, boxrule=1pt, pad at break*=1mm,colback=cellbackground, colframe=cellborder]
\prompt{In}{incolor}{111}{\boxspacing}
\begin{Verbatim}[commandchars=\\\{\}]
\PY{n}{plt}\PY{o}{.}\PY{n}{plot}\PY{p}{(}\PY{n}{xaxis\PYZus{}ne}\PY{p}{,} \PY{n}{ne\PYZus{}spectrum}\PY{p}{)}\PY{p}{;}
\end{Verbatim}
\end{tcolorbox}

    \begin{center}
    \adjustimage{max size={0.9\linewidth}{0.9\paperheight}}{output_57_0.png}
    \end{center}
    { \hspace*{\fill} \\}
    
    \textbf{\emph{Kriptón}}

    \begin{tcolorbox}[breakable, size=fbox, boxrule=1pt, pad at break*=1mm,colback=cellbackground, colframe=cellborder]
\prompt{In}{incolor}{112}{\boxspacing}
\begin{Verbatim}[commandchars=\\\{\}]
\PY{n}{plt}\PY{o}{.}\PY{n}{plot}\PY{p}{(}\PY{n}{xaxis\PYZus{}kr}\PY{p}{,} \PY{n}{kr\PYZus{}spectrum}\PY{p}{)}\PY{p}{;}
\end{Verbatim}
\end{tcolorbox}

    \begin{center}
    \adjustimage{max size={0.9\linewidth}{0.9\paperheight}}{output_59_0.png}
    \end{center}
    { \hspace*{\fill} \\}
    
    \textbf{\emph{Helio}}

    \begin{tcolorbox}[breakable, size=fbox, boxrule=1pt, pad at break*=1mm,colback=cellbackground, colframe=cellborder]
\prompt{In}{incolor}{113}{\boxspacing}
\begin{Verbatim}[commandchars=\\\{\}]
\PY{n}{plt}\PY{o}{.}\PY{n}{plot}\PY{p}{(}\PY{n}{xaxis\PYZus{}he}\PY{p}{,} \PY{n}{he\PYZus{}spectrum}\PY{p}{)}\PY{p}{;}
\end{Verbatim}
\end{tcolorbox}

    \begin{center}
    \adjustimage{max size={0.9\linewidth}{0.9\paperheight}}{output_61_0.png}
    \end{center}
    { \hspace*{\fill} \\}
    
    En este siguiente paso, tenemos que anotar qué valores (aproximados) de
píxel corresponden a longitudes de onda conocidas.

Adivinar esto simplemente mirando los espectros y comparándolos con
atlas espectrales es difícil y requiere experiencia.

El siguiente truco es averiguar en qué dirección aumenta la longitud de
onda: ¿hacia la izquierda o hacia la derecha? El mercurio no ayuda con
eso, ya que los dobletes azul-violeta y verde-amarillo están
aproximadamente a la misma distancia entre sí. Resulta que, a partir de
trabajo de prueba y error, la longitud de onda aumenta hacia la derecha.

En la siguiente línea de código lo que se hace es ubicar los pixeles que
se van a estudiar y las longitudes de onda teóricas que se tienen para
dichas líneas de emisión, tener presente que deben ser la misma cantidad
(las longitudes de onda se encuentran en nm)

    \textbf{\emph{Mercurio}}

    \begin{tcolorbox}[breakable, size=fbox, boxrule=1pt, pad at break*=1mm,colback=cellbackground, colframe=cellborder]
\prompt{In}{incolor}{221}{\boxspacing}
\begin{Verbatim}[commandchars=\\\{\}]
\PY{n}{guessed\PYZus{}wavelengths\PYZus{}hg} \PY{o}{=} \PY{p}{[}\PY{l+m+mf}{435.49}\PY{p}{,} \PY{l+m+mf}{491.23}\PY{p}{,} \PY{l+m+mf}{545.82}\PY{p}{,} \PY{l+m+mf}{567.65}\PY{p}{,} \PY{l+m+mf}{587.19}\PY{p}{,} \PY{l+m+mf}{606.73}\PY{p}{]}
\PY{n}{guessed\PYZus{}xvals\PYZus{}hg} \PY{o}{=} \PY{p}{[}\PY{l+m+mi}{100}\PY{p}{,} \PY{l+m+mi}{203}\PY{p}{,} \PY{l+m+mi}{302}\PY{p}{,} \PY{l+m+mi}{363}\PY{p}{,} \PY{l+m+mi}{408}\PY{p}{,} \PY{l+m+mi}{472}\PY{p}{]}
\end{Verbatim}
\end{tcolorbox}

    \hypertarget{mejorando-nuestras-suposiciones}{%
\subsection{Mejorando nuestras
suposiciones}\label{mejorando-nuestras-suposiciones}}

Podemos hacerlo mucho mejor al determinar los valores X de píxel tomando
la coordenada ponderada por la intensidad (momento 1):

    \begin{tcolorbox}[breakable, size=fbox, boxrule=1pt, pad at break*=1mm,colback=cellbackground, colframe=cellborder]
\prompt{In}{incolor}{222}{\boxspacing}
\begin{Verbatim}[commandchars=\\\{\}]
\PY{n}{npixels} \PY{o}{=} \PY{l+m+mi}{15}
\PY{n}{improved\PYZus{}xval\PYZus{}guesses\PYZus{}hg} \PY{o}{=} \PY{p}{[}\PY{n}{np}\PY{o}{.}\PY{n}{average}\PY{p}{(}\PY{n}{xaxis\PYZus{}hg}\PY{p}{[}\PY{n}{g}\PY{o}{\PYZhy{}}\PY{n}{npixels}\PY{p}{:}\PY{n}{g}\PY{o}{+}\PY{n}{npixels}\PY{p}{]}\PY{p}{,}
                                    \PY{n}{weights}\PY{o}{=}\PY{n}{hg\PYZus{}spectrum}\PY{p}{[}\PY{n}{g}\PY{o}{\PYZhy{}}\PY{n}{npixels}\PY{p}{:}\PY{n}{g}\PY{o}{+}\PY{n}{npixels}\PY{p}{]} \PY{o}{\PYZhy{}} \PY{n}{np}\PY{o}{.}\PY{n}{median}\PY{p}{(}\PY{n}{hg\PYZus{}spectrum}\PY{p}{)}\PY{p}{)}
                         \PY{k}{for} \PY{n}{g} \PY{o+ow}{in} \PY{n}{guessed\PYZus{}xvals\PYZus{}hg}\PY{p}{]}
\PY{n}{improved\PYZus{}xval\PYZus{}guesses\PYZus{}hg}
\end{Verbatim}
\end{tcolorbox}

            \begin{tcolorbox}[breakable, size=fbox, boxrule=.5pt, pad at break*=1mm, opacityfill=0]
\prompt{Out}{outcolor}{222}{\boxspacing}
\begin{Verbatim}[commandchars=\\\{\}]
[np.float64(101.01815345563979),
 np.float64(210.8963857840367),
 np.float64(304.20276711219304),
 np.float64(365.23598908219697),
 np.float64(409.5826508526847),
 np.float64(468.24336002098846)]
\end{Verbatim}
\end{tcolorbox}
        
    Observamos los valores de nuestras supocisiones y los mejorados, si
coiciden con los picos seleccionados se está realizando de forma óptima,
si quedan muy lejos hay que cambiar el valor del pixel de nuestras
suposiciones

    \begin{tcolorbox}[breakable, size=fbox, boxrule=1pt, pad at break*=1mm,colback=cellbackground, colframe=cellborder]
\prompt{In}{incolor}{223}{\boxspacing}
\begin{Verbatim}[commandchars=\\\{\}]
\PY{n}{plt}\PY{o}{.}\PY{n}{plot}\PY{p}{(}\PY{n}{xaxis\PYZus{}hg}\PY{p}{,} \PY{n}{hg\PYZus{}spectrum}\PY{p}{,} \PY{n}{label}\PY{o}{=}\PY{l+s+s1}{\PYZsq{}}\PY{l+s+s1}{Espectro Hg}\PY{l+s+s1}{\PYZsq{}}\PY{p}{)}
\PY{n}{plt}\PY{o}{.}\PY{n}{plot}\PY{p}{(}\PY{n}{guessed\PYZus{}xvals\PYZus{}hg}\PY{p}{,} \PY{p}{[}\PY{l+m+mf}{0.5}\PY{p}{]}\PY{o}{*}\PY{l+m+mi}{6}\PY{p}{,} \PY{l+s+s1}{\PYZsq{}}\PY{l+s+s1}{rx}\PY{l+s+s1}{\PYZsq{}}\PY{p}{,} \PY{n}{label}\PY{o}{=}\PY{l+s+s1}{\PYZsq{}}\PY{l+s+s1}{Estimación inicial}\PY{l+s+s1}{\PYZsq{}}\PY{p}{)}
\PY{n}{plt}\PY{o}{.}\PY{n}{plot}\PY{p}{(}\PY{n}{improved\PYZus{}xval\PYZus{}guesses\PYZus{}hg}\PY{p}{,} \PY{p}{[}\PY{l+m+mf}{0.5}\PY{p}{]}\PY{o}{*}\PY{l+m+mi}{6}\PY{p}{,} \PY{l+s+s1}{\PYZsq{}}\PY{l+s+s1}{g+}\PY{l+s+s1}{\PYZsq{}}\PY{p}{,} \PY{n}{label}\PY{o}{=}\PY{l+s+s1}{\PYZsq{}}\PY{l+s+s1}{Centro refinado}\PY{l+s+s1}{\PYZsq{}}\PY{p}{)}
\PY{n}{plt}\PY{o}{.}\PY{n}{legend}\PY{p}{(}\PY{p}{)}
\PY{n}{plt}\PY{o}{.}\PY{n}{xlabel}\PY{p}{(}\PY{l+s+s1}{\PYZsq{}}\PY{l+s+s1}{Píxel}\PY{l+s+s1}{\PYZsq{}}\PY{p}{)}
\PY{n}{plt}\PY{o}{.}\PY{n}{ylabel}\PY{p}{(}\PY{l+s+s1}{\PYZsq{}}\PY{l+s+s1}{Intensidad}\PY{l+s+s1}{\PYZsq{}}\PY{p}{)}
\PY{n}{plt}\PY{o}{.}\PY{n}{show}\PY{p}{(}\PY{p}{)}
\end{Verbatim}
\end{tcolorbox}

    \begin{center}
    \adjustimage{max size={0.9\linewidth}{0.9\paperheight}}{output_68_0.png}
    \end{center}
    { \hspace*{\fill} \\}
    
    \textbf{\emph{Neón}}

    \begin{tcolorbox}[breakable, size=fbox, boxrule=1pt, pad at break*=1mm,colback=cellbackground, colframe=cellborder]
\prompt{In}{incolor}{153}{\boxspacing}
\begin{Verbatim}[commandchars=\\\{\}]
\PY{n}{guessed\PYZus{}wavelengths\PYZus{}ne} \PY{o}{=} \PY{p}{[}\PY{l+m+mf}{584.47}\PY{p}{,} \PY{l+m+mf}{587.36}\PY{p}{,} \PY{l+m+mf}{594.66}\PY{p}{,} \PY{l+m+mf}{597.53}\PY{p}{,} 
                          \PY{l+m+mf}{603.28}\PY{p}{,} \PY{l+m+mf}{607.3}\PY{p}{,} \PY{l+m+mf}{617.65}\PY{p}{,} \PY{l+m+mf}{622.24}\PY{p}{,} \PY{l+m+mf}{626.84}\PY{p}{]}
\PY{n}{guessed\PYZus{}xvals\PYZus{}ne} \PY{o}{=} \PY{p}{[}\PY{l+m+mi}{130}\PY{p}{,} \PY{l+m+mi}{190}\PY{p}{,} \PY{l+m+mi}{254}\PY{p}{,} \PY{l+m+mi}{290}\PY{p}{,} \PY{l+m+mi}{359}\PY{p}{,} \PY{l+m+mi}{390}\PY{p}{,} \PY{l+m+mi}{418}\PY{p}{,} \PY{l+m+mi}{492}\PY{p}{,} \PY{l+m+mi}{588}\PY{p}{]}
\end{Verbatim}
\end{tcolorbox}

    \begin{tcolorbox}[breakable, size=fbox, boxrule=1pt, pad at break*=1mm,colback=cellbackground, colframe=cellborder]
\prompt{In}{incolor}{154}{\boxspacing}
\begin{Verbatim}[commandchars=\\\{\}]
\PY{n}{npixels} \PY{o}{=} \PY{l+m+mi}{15}
\PY{n}{improved\PYZus{}xval\PYZus{}guesses\PYZus{}ne} \PY{o}{=} \PY{p}{[}\PY{n}{np}\PY{o}{.}\PY{n}{average}\PY{p}{(}\PY{n}{xaxis\PYZus{}ne}\PY{p}{[}\PY{n}{g}\PY{o}{\PYZhy{}}\PY{n}{npixels}\PY{p}{:}\PY{n}{g}\PY{o}{+}\PY{n}{npixels}\PY{p}{]}\PY{p}{,}
                                    \PY{n}{weights}\PY{o}{=}\PY{n}{ne\PYZus{}spectrum}\PY{p}{[}\PY{n}{g}\PY{o}{\PYZhy{}}\PY{n}{npixels}\PY{p}{:}\PY{n}{g}\PY{o}{+}\PY{n}{npixels}\PY{p}{]} \PY{o}{\PYZhy{}} \PY{n}{np}\PY{o}{.}\PY{n}{median}\PY{p}{(}\PY{n}{ne\PYZus{}spectrum}\PY{p}{)}\PY{p}{)}
                         \PY{k}{for} \PY{n}{g} \PY{o+ow}{in} \PY{n}{guessed\PYZus{}xvals\PYZus{}ne}\PY{p}{]}
\PY{n}{improved\PYZus{}xval\PYZus{}guesses\PYZus{}ne}
\end{Verbatim}
\end{tcolorbox}

            \begin{tcolorbox}[breakable, size=fbox, boxrule=.5pt, pad at break*=1mm, opacityfill=0]
\prompt{Out}{outcolor}{154}{\boxspacing}
\begin{Verbatim}[commandchars=\\\{\}]
[np.float64(133.01102911965515),
 np.float64(185.5284212469613),
 np.float64(254.7748769737537),
 np.float64(288.41034410972594),
 np.float64(356.4902213968072),
 np.float64(388.7223531978048),
 np.float64(418.2268099738853),
 np.float64(486.88728194468274),
 np.float64(582.5898134276622)]
\end{Verbatim}
\end{tcolorbox}
        
    \begin{tcolorbox}[breakable, size=fbox, boxrule=1pt, pad at break*=1mm,colback=cellbackground, colframe=cellborder]
\prompt{In}{incolor}{155}{\boxspacing}
\begin{Verbatim}[commandchars=\\\{\}]
\PY{n}{plt}\PY{o}{.}\PY{n}{plot}\PY{p}{(}\PY{n}{xaxis\PYZus{}ne}\PY{p}{,} \PY{n}{ne\PYZus{}spectrum}\PY{p}{,} \PY{n}{label}\PY{o}{=}\PY{l+s+s1}{\PYZsq{}}\PY{l+s+s1}{Espectro Neon}\PY{l+s+s1}{\PYZsq{}}\PY{p}{)}
\PY{n}{plt}\PY{o}{.}\PY{n}{plot}\PY{p}{(}\PY{n}{guessed\PYZus{}xvals\PYZus{}ne}\PY{p}{,} \PY{p}{[}\PY{l+m+mf}{0.5}\PY{p}{]}\PY{o}{*}\PY{l+m+mi}{9}\PY{p}{,} \PY{l+s+s1}{\PYZsq{}}\PY{l+s+s1}{rx}\PY{l+s+s1}{\PYZsq{}}\PY{p}{,} \PY{n}{label}\PY{o}{=}\PY{l+s+s1}{\PYZsq{}}\PY{l+s+s1}{Estimación inicial}\PY{l+s+s1}{\PYZsq{}}\PY{p}{)}
\PY{n}{plt}\PY{o}{.}\PY{n}{plot}\PY{p}{(}\PY{n}{improved\PYZus{}xval\PYZus{}guesses\PYZus{}ne}\PY{p}{,} \PY{p}{[}\PY{l+m+mf}{0.5}\PY{p}{]}\PY{o}{*}\PY{l+m+mi}{9}\PY{p}{,} \PY{l+s+s1}{\PYZsq{}}\PY{l+s+s1}{g+}\PY{l+s+s1}{\PYZsq{}}\PY{p}{,} \PY{n}{label}\PY{o}{=}\PY{l+s+s1}{\PYZsq{}}\PY{l+s+s1}{Centro refinado}\PY{l+s+s1}{\PYZsq{}}\PY{p}{)}
\PY{n}{plt}\PY{o}{.}\PY{n}{legend}\PY{p}{(}\PY{p}{)}
\PY{n}{plt}\PY{o}{.}\PY{n}{xlabel}\PY{p}{(}\PY{l+s+s1}{\PYZsq{}}\PY{l+s+s1}{Píxel}\PY{l+s+s1}{\PYZsq{}}\PY{p}{)}
\PY{n}{plt}\PY{o}{.}\PY{n}{ylabel}\PY{p}{(}\PY{l+s+s1}{\PYZsq{}}\PY{l+s+s1}{Intensidad}\PY{l+s+s1}{\PYZsq{}}\PY{p}{)}
\PY{n}{plt}\PY{o}{.}\PY{n}{show}\PY{p}{(}\PY{p}{)}
\end{Verbatim}
\end{tcolorbox}

    \begin{center}
    \adjustimage{max size={0.9\linewidth}{0.9\paperheight}}{output_72_0.png}
    \end{center}
    { \hspace*{\fill} \\}
    
    \textbf{\emph{Krptón}}

    \begin{tcolorbox}[breakable, size=fbox, boxrule=1pt, pad at break*=1mm,colback=cellbackground, colframe=cellborder]
\prompt{In}{incolor}{179}{\boxspacing}
\begin{Verbatim}[commandchars=\\\{\}]
\PY{n}{guessed\PYZus{}wavelengths\PYZus{}kr} \PY{o}{=} \PY{p}{[}\PY{l+m+mf}{556.73}\PY{p}{,} \PY{l+m+mf}{587.19}\PY{p}{,} \PY{l+m+mf}{605.58}\PY{p}{,} \PY{l+m+mf}{642.36}\PY{p}{,} \PY{l+m+mf}{645.8}\PY{p}{]}
\PY{n}{guessed\PYZus{}xvals\PYZus{}kr} \PY{o}{=} \PY{p}{[}\PY{l+m+mi}{280}\PY{p}{,} \PY{l+m+mi}{352}\PY{p}{,} \PY{l+m+mi}{392}\PY{p}{,} \PY{l+m+mi}{484}\PY{p}{,} \PY{l+m+mi}{509}\PY{p}{]}
\end{Verbatim}
\end{tcolorbox}

    \begin{tcolorbox}[breakable, size=fbox, boxrule=1pt, pad at break*=1mm,colback=cellbackground, colframe=cellborder]
\prompt{In}{incolor}{180}{\boxspacing}
\begin{Verbatim}[commandchars=\\\{\}]
\PY{n}{npixels} \PY{o}{=} \PY{l+m+mi}{15}
\PY{n}{improved\PYZus{}xval\PYZus{}guesses\PYZus{}kr} \PY{o}{=} \PY{p}{[}\PY{n}{np}\PY{o}{.}\PY{n}{average}\PY{p}{(}\PY{n}{xaxis\PYZus{}kr}\PY{p}{[}\PY{n}{g}\PY{o}{\PYZhy{}}\PY{n}{npixels}\PY{p}{:}\PY{n}{g}\PY{o}{+}\PY{n}{npixels}\PY{p}{]}\PY{p}{,}
                                    \PY{n}{weights}\PY{o}{=}\PY{n}{kr\PYZus{}spectrum}\PY{p}{[}\PY{n}{g}\PY{o}{\PYZhy{}}\PY{n}{npixels}\PY{p}{:}\PY{n}{g}\PY{o}{+}\PY{n}{npixels}\PY{p}{]} \PY{o}{\PYZhy{}} \PY{n}{np}\PY{o}{.}\PY{n}{median}\PY{p}{(}\PY{n}{kr\PYZus{}spectrum}\PY{p}{)}\PY{p}{)}
                         \PY{k}{for} \PY{n}{g} \PY{o+ow}{in} \PY{n}{guessed\PYZus{}xvals\PYZus{}kr}\PY{p}{]}
\PY{n}{improved\PYZus{}xval\PYZus{}guesses\PYZus{}kr}
\end{Verbatim}
\end{tcolorbox}

            \begin{tcolorbox}[breakable, size=fbox, boxrule=.5pt, pad at break*=1mm, opacityfill=0]
\prompt{Out}{outcolor}{180}{\boxspacing}
\begin{Verbatim}[commandchars=\\\{\}]
[np.float64(287.53134824826685),
 np.float64(354.8258707526527),
 np.float64(392.70848061769647),
 np.float64(482.87929801092395),
 np.float64(509.05712702013045)]
\end{Verbatim}
\end{tcolorbox}
        
    \begin{tcolorbox}[breakable, size=fbox, boxrule=1pt, pad at break*=1mm,colback=cellbackground, colframe=cellborder]
\prompt{In}{incolor}{181}{\boxspacing}
\begin{Verbatim}[commandchars=\\\{\}]
\PY{n}{plt}\PY{o}{.}\PY{n}{plot}\PY{p}{(}\PY{n}{xaxis\PYZus{}kr}\PY{p}{,} \PY{n}{kr\PYZus{}spectrum}\PY{p}{,} \PY{n}{label}\PY{o}{=}\PY{l+s+s1}{\PYZsq{}}\PY{l+s+s1}{Espectro Kriptón}\PY{l+s+s1}{\PYZsq{}}\PY{p}{)}
\PY{n}{plt}\PY{o}{.}\PY{n}{plot}\PY{p}{(}\PY{n}{guessed\PYZus{}xvals\PYZus{}kr}\PY{p}{,} \PY{p}{[}\PY{l+m+mf}{0.2}\PY{p}{]}\PY{o}{*}\PY{l+m+mi}{5}\PY{p}{,} \PY{l+s+s1}{\PYZsq{}}\PY{l+s+s1}{rx}\PY{l+s+s1}{\PYZsq{}}\PY{p}{,} \PY{n}{label}\PY{o}{=}\PY{l+s+s1}{\PYZsq{}}\PY{l+s+s1}{Estimación inicial}\PY{l+s+s1}{\PYZsq{}}\PY{p}{)}
\PY{n}{plt}\PY{o}{.}\PY{n}{plot}\PY{p}{(}\PY{n}{improved\PYZus{}xval\PYZus{}guesses\PYZus{}kr}\PY{p}{,} \PY{p}{[}\PY{l+m+mf}{0.2}\PY{p}{]}\PY{o}{*}\PY{l+m+mi}{5}\PY{p}{,} \PY{l+s+s1}{\PYZsq{}}\PY{l+s+s1}{g+}\PY{l+s+s1}{\PYZsq{}}\PY{p}{,} \PY{n}{label}\PY{o}{=}\PY{l+s+s1}{\PYZsq{}}\PY{l+s+s1}{Centro refinado}\PY{l+s+s1}{\PYZsq{}}\PY{p}{)}
\PY{n}{plt}\PY{o}{.}\PY{n}{legend}\PY{p}{(}\PY{p}{)}
\PY{n}{plt}\PY{o}{.}\PY{n}{xlabel}\PY{p}{(}\PY{l+s+s1}{\PYZsq{}}\PY{l+s+s1}{Píxel}\PY{l+s+s1}{\PYZsq{}}\PY{p}{)}
\PY{n}{plt}\PY{o}{.}\PY{n}{ylabel}\PY{p}{(}\PY{l+s+s1}{\PYZsq{}}\PY{l+s+s1}{Intensidad}\PY{l+s+s1}{\PYZsq{}}\PY{p}{)}
\PY{n}{plt}\PY{o}{.}\PY{n}{show}\PY{p}{(}\PY{p}{)}
\end{Verbatim}
\end{tcolorbox}

    \begin{center}
    \adjustimage{max size={0.9\linewidth}{0.9\paperheight}}{output_76_0.png}
    \end{center}
    { \hspace*{\fill} \\}
    
    \textbf{\emph{Helio}}

    \begin{tcolorbox}[breakable, size=fbox, boxrule=1pt, pad at break*=1mm,colback=cellbackground, colframe=cellborder]
\prompt{In}{incolor}{199}{\boxspacing}
\begin{Verbatim}[commandchars=\\\{\}]
\PY{n}{guessed\PYZus{}wavelengths\PYZus{}he} \PY{o}{=} \PY{p}{[}\PY{l+m+mf}{446.98}\PY{p}{,} \PY{l+m+mf}{491.8}\PY{p}{,} \PY{l+m+mf}{587.19}\PY{p}{,} \PY{l+m+mf}{668.21}\PY{p}{]}
\PY{n}{guessed\PYZus{}xvals\PYZus{}he} \PY{o}{=} \PY{p}{[}\PY{l+m+mi}{140}\PY{p}{,} \PY{l+m+mi}{270}\PY{p}{,} \PY{l+m+mi}{470}\PY{p}{,} \PY{l+m+mi}{660}\PY{p}{]}
\end{Verbatim}
\end{tcolorbox}

    \begin{tcolorbox}[breakable, size=fbox, boxrule=1pt, pad at break*=1mm,colback=cellbackground, colframe=cellborder]
\prompt{In}{incolor}{200}{\boxspacing}
\begin{Verbatim}[commandchars=\\\{\}]
\PY{n}{npixels} \PY{o}{=} \PY{l+m+mi}{15}
\PY{n}{improved\PYZus{}xval\PYZus{}guesses\PYZus{}he} \PY{o}{=} \PY{p}{[}\PY{n}{np}\PY{o}{.}\PY{n}{average}\PY{p}{(}\PY{n}{xaxis\PYZus{}he}\PY{p}{[}\PY{n}{g}\PY{o}{\PYZhy{}}\PY{n}{npixels}\PY{p}{:}\PY{n}{g}\PY{o}{+}\PY{n}{npixels}\PY{p}{]}\PY{p}{,}
                                    \PY{n}{weights}\PY{o}{=}\PY{n}{he\PYZus{}spectrum}\PY{p}{[}\PY{n}{g}\PY{o}{\PYZhy{}}\PY{n}{npixels}\PY{p}{:}\PY{n}{g}\PY{o}{+}\PY{n}{npixels}\PY{p}{]} \PY{o}{\PYZhy{}} \PY{n}{np}\PY{o}{.}\PY{n}{median}\PY{p}{(}\PY{n}{he\PYZus{}spectrum}\PY{p}{)}\PY{p}{)}
                         \PY{k}{for} \PY{n}{g} \PY{o+ow}{in} \PY{n}{guessed\PYZus{}xvals\PYZus{}he}\PY{p}{]}
\PY{n}{improved\PYZus{}xval\PYZus{}guesses\PYZus{}he}
\end{Verbatim}
\end{tcolorbox}

            \begin{tcolorbox}[breakable, size=fbox, boxrule=.5pt, pad at break*=1mm, opacityfill=0]
\prompt{Out}{outcolor}{200}{\boxspacing}
\begin{Verbatim}[commandchars=\\\{\}]
[np.float64(137.6654284329619),
 np.float64(266.2029225364427),
 np.float64(470.3120139965751),
 np.float64(655.9549174964578)]
\end{Verbatim}
\end{tcolorbox}
        
    \begin{tcolorbox}[breakable, size=fbox, boxrule=1pt, pad at break*=1mm,colback=cellbackground, colframe=cellborder]
\prompt{In}{incolor}{201}{\boxspacing}
\begin{Verbatim}[commandchars=\\\{\}]
\PY{n}{plt}\PY{o}{.}\PY{n}{plot}\PY{p}{(}\PY{n}{xaxis\PYZus{}he}\PY{p}{,} \PY{n}{he\PYZus{}spectrum}\PY{p}{,} \PY{n}{label}\PY{o}{=}\PY{l+s+s1}{\PYZsq{}}\PY{l+s+s1}{Espectro Helio}\PY{l+s+s1}{\PYZsq{}}\PY{p}{)}
\PY{n}{plt}\PY{o}{.}\PY{n}{plot}\PY{p}{(}\PY{n}{guessed\PYZus{}xvals\PYZus{}he}\PY{p}{,} \PY{p}{[}\PY{l+m+mf}{0.7}\PY{p}{]}\PY{o}{*}\PY{l+m+mi}{4}\PY{p}{,} \PY{l+s+s1}{\PYZsq{}}\PY{l+s+s1}{rx}\PY{l+s+s1}{\PYZsq{}}\PY{p}{,} \PY{n}{label}\PY{o}{=}\PY{l+s+s1}{\PYZsq{}}\PY{l+s+s1}{Estimación inicial}\PY{l+s+s1}{\PYZsq{}}\PY{p}{)}
\PY{n}{plt}\PY{o}{.}\PY{n}{plot}\PY{p}{(}\PY{n}{improved\PYZus{}xval\PYZus{}guesses\PYZus{}he}\PY{p}{,} \PY{p}{[}\PY{l+m+mf}{0.7}\PY{p}{]}\PY{o}{*}\PY{l+m+mi}{4}\PY{p}{,} \PY{l+s+s1}{\PYZsq{}}\PY{l+s+s1}{g+}\PY{l+s+s1}{\PYZsq{}}\PY{p}{,} \PY{n}{label}\PY{o}{=}\PY{l+s+s1}{\PYZsq{}}\PY{l+s+s1}{Centro refinado}\PY{l+s+s1}{\PYZsq{}}\PY{p}{)}
\PY{n}{plt}\PY{o}{.}\PY{n}{legend}\PY{p}{(}\PY{p}{)}
\PY{n}{plt}\PY{o}{.}\PY{n}{xlabel}\PY{p}{(}\PY{l+s+s1}{\PYZsq{}}\PY{l+s+s1}{Píxel}\PY{l+s+s1}{\PYZsq{}}\PY{p}{)}
\PY{n}{plt}\PY{o}{.}\PY{n}{ylabel}\PY{p}{(}\PY{l+s+s1}{\PYZsq{}}\PY{l+s+s1}{Intensidad}\PY{l+s+s1}{\PYZsq{}}\PY{p}{)}
\PY{n}{plt}\PY{o}{.}\PY{n}{show}\PY{p}{(}\PY{p}{)}
\end{Verbatim}
\end{tcolorbox}

    \begin{center}
    \adjustimage{max size={0.9\linewidth}{0.9\paperheight}}{output_80_0.png}
    \end{center}
    { \hspace*{\fill} \\}
    
    \begin{tcolorbox}[breakable, size=fbox, boxrule=1pt, pad at break*=1mm,colback=cellbackground, colframe=cellborder]
\prompt{In}{incolor}{224}{\boxspacing}
\begin{Verbatim}[commandchars=\\\{\}]
\PY{n}{linfitter} \PY{o}{=} \PY{n}{LinearLSQFitter}\PY{p}{(}\PY{p}{)}
\end{Verbatim}
\end{tcolorbox}

    Usamos un modelo \texttt{Linear1D} porque más adelante querremos usar su
inversa (otros modelos no son invertibles).

    Se grafica ahora en longitudes de onda, en vez de pixeles

    \textbf{\emph{Mercurio}}

    \begin{tcolorbox}[breakable, size=fbox, boxrule=1pt, pad at break*=1mm,colback=cellbackground, colframe=cellborder]
\prompt{In}{incolor}{231}{\boxspacing}
\begin{Verbatim}[commandchars=\\\{\}]
\PY{n}{wlmodel} \PY{o}{=} \PY{n}{Linear1D}\PY{p}{(}\PY{p}{)}
\PY{n}{linfit\PYZus{}wlmodel\PYZus{}hg} \PY{o}{=} \PY{n}{linfitter}\PY{p}{(}\PY{n}{model}\PY{o}{=}\PY{n}{wlmodel}\PY{p}{,} \PY{n}{x}\PY{o}{=}\PY{n}{improved\PYZus{}xval\PYZus{}guesses\PYZus{}hg}\PY{p}{,} \PY{n}{y}\PY{o}{=}\PY{n}{guessed\PYZus{}wavelengths\PYZus{}hg}\PY{p}{)}
\PY{n}{wavelengths\PYZus{}hg} \PY{o}{=} \PY{n}{linfit\PYZus{}wlmodel\PYZus{}hg}\PY{p}{(}\PY{n}{xaxis\PYZus{}hg}\PY{p}{)} \PY{o}{*} \PY{n}{u}\PY{o}{.}\PY{n}{nm}
\PY{n}{linfit\PYZus{}wlmodel\PYZus{}hg}
\end{Verbatim}
\end{tcolorbox}

            \begin{tcolorbox}[breakable, size=fbox, boxrule=.5pt, pad at break*=1mm, opacityfill=0]
\prompt{Out}{outcolor}{231}{\boxspacing}
\begin{Verbatim}[commandchars=\\\{\}]
<Linear1D(slope=0.47440419, intercept=392.01792493)>
\end{Verbatim}
\end{tcolorbox}
        
    \begin{tcolorbox}[breakable, size=fbox, boxrule=1pt, pad at break*=1mm,colback=cellbackground, colframe=cellborder]
\prompt{In}{incolor}{232}{\boxspacing}
\begin{Verbatim}[commandchars=\\\{\}]
\PY{n}{plt}\PY{o}{.}\PY{n}{plot}\PY{p}{(}\PY{n}{wavelengths\PYZus{}hg}\PY{p}{,} \PY{n}{hg\PYZus{}spectrum}\PY{p}{)}
\PY{n}{plt}\PY{o}{.}\PY{n}{plot}\PY{p}{(}\PY{n}{guessed\PYZus{}wavelengths\PYZus{}hg}\PY{p}{,} \PY{p}{[}\PY{l+m+mf}{0.7}\PY{p}{]}\PY{o}{*}\PY{l+m+mi}{6}\PY{p}{,} \PY{l+s+s1}{\PYZsq{}}\PY{l+s+s1}{x}\PY{l+s+s1}{\PYZsq{}}\PY{p}{)}\PY{p}{;}
\end{Verbatim}
\end{tcolorbox}

    \begin{center}
    \adjustimage{max size={0.9\linewidth}{0.9\paperheight}}{output_86_0.png}
    \end{center}
    { \hspace*{\fill} \\}
    
    \textbf{\emph{Neón}}

    \begin{tcolorbox}[breakable, size=fbox, boxrule=1pt, pad at break*=1mm,colback=cellbackground, colframe=cellborder]
\prompt{In}{incolor}{234}{\boxspacing}
\begin{Verbatim}[commandchars=\\\{\}]
\PY{n}{wlmodel} \PY{o}{=} \PY{n}{Linear1D}\PY{p}{(}\PY{p}{)}
\PY{n}{linfit\PYZus{}wlmodel\PYZus{}ne} \PY{o}{=} \PY{n}{linfitter}\PY{p}{(}\PY{n}{model}\PY{o}{=}\PY{n}{wlmodel}\PY{p}{,} \PY{n}{x}\PY{o}{=}\PY{n}{improved\PYZus{}xval\PYZus{}guesses\PYZus{}ne}\PY{p}{,} \PY{n}{y}\PY{o}{=}\PY{n}{guessed\PYZus{}wavelengths\PYZus{}ne}\PY{p}{)}
\PY{n}{wavelengths\PYZus{}ne} \PY{o}{=} \PY{n}{linfit\PYZus{}wlmodel\PYZus{}ne}\PY{p}{(}\PY{n}{xaxis\PYZus{}ne}\PY{p}{)} \PY{o}{*} \PY{n}{u}\PY{o}{.}\PY{n}{nm}
\PY{n}{linfit\PYZus{}wlmodel\PYZus{}ne}
\end{Verbatim}
\end{tcolorbox}

            \begin{tcolorbox}[breakable, size=fbox, boxrule=.5pt, pad at break*=1mm, opacityfill=0]
\prompt{Out}{outcolor}{234}{\boxspacing}
\begin{Verbatim}[commandchars=\\\{\}]
<Linear1D(slope=0.10356106, intercept=568.98285354)>
\end{Verbatim}
\end{tcolorbox}
        
    \begin{tcolorbox}[breakable, size=fbox, boxrule=1pt, pad at break*=1mm,colback=cellbackground, colframe=cellborder]
\prompt{In}{incolor}{235}{\boxspacing}
\begin{Verbatim}[commandchars=\\\{\}]
\PY{n}{plt}\PY{o}{.}\PY{n}{plot}\PY{p}{(}\PY{n}{wavelengths\PYZus{}ne}\PY{p}{,} \PY{n}{ne\PYZus{}spectrum}\PY{p}{)}
\PY{n}{plt}\PY{o}{.}\PY{n}{plot}\PY{p}{(}\PY{n}{guessed\PYZus{}wavelengths\PYZus{}ne}\PY{p}{,} \PY{p}{[}\PY{l+m+mf}{0.5}\PY{p}{]}\PY{o}{*}\PY{l+m+mi}{9}\PY{p}{,} \PY{l+s+s1}{\PYZsq{}}\PY{l+s+s1}{x}\PY{l+s+s1}{\PYZsq{}}\PY{p}{)}\PY{p}{;}
\end{Verbatim}
\end{tcolorbox}

    \begin{center}
    \adjustimage{max size={0.9\linewidth}{0.9\paperheight}}{output_89_0.png}
    \end{center}
    { \hspace*{\fill} \\}
    
    \textbf{\emph{Kriptón}}

    \begin{tcolorbox}[breakable, size=fbox, boxrule=1pt, pad at break*=1mm,colback=cellbackground, colframe=cellborder]
\prompt{In}{incolor}{236}{\boxspacing}
\begin{Verbatim}[commandchars=\\\{\}]
\PY{n}{wlmodel} \PY{o}{=} \PY{n}{Linear1D}\PY{p}{(}\PY{p}{)}
\PY{n}{linfit\PYZus{}wlmodel\PYZus{}kr} \PY{o}{=} \PY{n}{linfitter}\PY{p}{(}\PY{n}{model}\PY{o}{=}\PY{n}{wlmodel}\PY{p}{,} \PY{n}{x}\PY{o}{=}\PY{n}{improved\PYZus{}xval\PYZus{}guesses\PYZus{}kr}\PY{p}{,} \PY{n}{y}\PY{o}{=}\PY{n}{guessed\PYZus{}wavelengths\PYZus{}kr}\PY{p}{)}
\PY{n}{wavelengths\PYZus{}kr} \PY{o}{=} \PY{n}{linfit\PYZus{}wlmodel\PYZus{}kr}\PY{p}{(}\PY{n}{xaxis\PYZus{}kr}\PY{p}{)} \PY{o}{*} \PY{n}{u}\PY{o}{.}\PY{n}{nm}
\PY{n}{linfit\PYZus{}wlmodel\PYZus{}kr}
\end{Verbatim}
\end{tcolorbox}

            \begin{tcolorbox}[breakable, size=fbox, boxrule=.5pt, pad at break*=1mm, opacityfill=0]
\prompt{Out}{outcolor}{236}{\boxspacing}
\begin{Verbatim}[commandchars=\\\{\}]
<Linear1D(slope=0.41087675, intercept=440.96239182)>
\end{Verbatim}
\end{tcolorbox}
        
    \begin{tcolorbox}[breakable, size=fbox, boxrule=1pt, pad at break*=1mm,colback=cellbackground, colframe=cellborder]
\prompt{In}{incolor}{238}{\boxspacing}
\begin{Verbatim}[commandchars=\\\{\}]
\PY{n}{plt}\PY{o}{.}\PY{n}{plot}\PY{p}{(}\PY{n}{wavelengths\PYZus{}kr}\PY{p}{,} \PY{n}{kr\PYZus{}spectrum}\PY{p}{)}
\PY{n}{plt}\PY{o}{.}\PY{n}{plot}\PY{p}{(}\PY{n}{guessed\PYZus{}wavelengths\PYZus{}kr}\PY{p}{,} \PY{p}{[}\PY{l+m+mf}{0.2}\PY{p}{]}\PY{o}{*}\PY{l+m+mi}{5}\PY{p}{,} \PY{l+s+s1}{\PYZsq{}}\PY{l+s+s1}{x}\PY{l+s+s1}{\PYZsq{}}\PY{p}{)}\PY{p}{;}
\end{Verbatim}
\end{tcolorbox}

    \begin{center}
    \adjustimage{max size={0.9\linewidth}{0.9\paperheight}}{output_92_0.png}
    \end{center}
    { \hspace*{\fill} \\}
    
    \textbf{\emph{Helio}}

    \begin{tcolorbox}[breakable, size=fbox, boxrule=1pt, pad at break*=1mm,colback=cellbackground, colframe=cellborder]
\prompt{In}{incolor}{239}{\boxspacing}
\begin{Verbatim}[commandchars=\\\{\}]
\PY{n}{wlmodel} \PY{o}{=} \PY{n}{Linear1D}\PY{p}{(}\PY{p}{)}
\PY{n}{linfit\PYZus{}wlmodel\PYZus{}he} \PY{o}{=} \PY{n}{linfitter}\PY{p}{(}\PY{n}{model}\PY{o}{=}\PY{n}{wlmodel}\PY{p}{,} \PY{n}{x}\PY{o}{=}\PY{n}{improved\PYZus{}xval\PYZus{}guesses\PYZus{}he}\PY{p}{,} \PY{n}{y}\PY{o}{=}\PY{n}{guessed\PYZus{}wavelengths\PYZus{}he}\PY{p}{)}
\PY{n}{wavelengths\PYZus{}he} \PY{o}{=} \PY{n}{linfit\PYZus{}wlmodel\PYZus{}he}\PY{p}{(}\PY{n}{xaxis\PYZus{}he}\PY{p}{)} \PY{o}{*} \PY{n}{u}\PY{o}{.}\PY{n}{nm}
\PY{n}{linfit\PYZus{}wlmodel\PYZus{}he}
\end{Verbatim}
\end{tcolorbox}

            \begin{tcolorbox}[breakable, size=fbox, boxrule=.5pt, pad at break*=1mm, opacityfill=0]
\prompt{Out}{outcolor}{239}{\boxspacing}
\begin{Verbatim}[commandchars=\\\{\}]
<Linear1D(slope=0.43333839, intercept=382.77840874)>
\end{Verbatim}
\end{tcolorbox}
        
    \begin{tcolorbox}[breakable, size=fbox, boxrule=1pt, pad at break*=1mm,colback=cellbackground, colframe=cellborder]
\prompt{In}{incolor}{241}{\boxspacing}
\begin{Verbatim}[commandchars=\\\{\}]
\PY{n}{plt}\PY{o}{.}\PY{n}{plot}\PY{p}{(}\PY{n}{wavelengths\PYZus{}he}\PY{p}{,} \PY{n}{he\PYZus{}spectrum}\PY{p}{)}
\PY{n}{plt}\PY{o}{.}\PY{n}{plot}\PY{p}{(}\PY{n}{guessed\PYZus{}wavelengths\PYZus{}he}\PY{p}{,} \PY{p}{[}\PY{l+m+mf}{0.7}\PY{p}{]}\PY{o}{*}\PY{l+m+mi}{4}\PY{p}{,} \PY{l+s+s1}{\PYZsq{}}\PY{l+s+s1}{x}\PY{l+s+s1}{\PYZsq{}}\PY{p}{)}\PY{p}{;}
\end{Verbatim}
\end{tcolorbox}

    \begin{center}
    \adjustimage{max size={0.9\linewidth}{0.9\paperheight}}{output_95_0.png}
    \end{center}
    { \hspace*{\fill} \\}
    
    Mostramos nuestro modelo \(\lambda(x)\) frente a las ``suposiciones'' de
entrada:

    \textbf{\emph{Mercurio}}

    \begin{tcolorbox}[breakable, size=fbox, boxrule=1pt, pad at break*=1mm,colback=cellbackground, colframe=cellborder]
\prompt{In}{incolor}{242}{\boxspacing}
\begin{Verbatim}[commandchars=\\\{\}]
\PY{n}{plt}\PY{o}{.}\PY{n}{plot}\PY{p}{(}\PY{n}{improved\PYZus{}xval\PYZus{}guesses\PYZus{}hg}\PY{p}{,} \PY{n}{guessed\PYZus{}wavelengths\PYZus{}hg}\PY{p}{,} \PY{l+s+s1}{\PYZsq{}}\PY{l+s+s1}{o}\PY{l+s+s1}{\PYZsq{}}\PY{p}{)}
\PY{n}{plt}\PY{o}{.}\PY{n}{plot}\PY{p}{(}\PY{n}{xaxis\PYZus{}hg}\PY{p}{,} \PY{n}{wavelengths\PYZus{}hg}\PY{p}{,} \PY{l+s+s1}{\PYZsq{}}\PY{l+s+s1}{\PYZhy{}}\PY{l+s+s1}{\PYZsq{}}\PY{p}{)}
\PY{n}{plt}\PY{o}{.}\PY{n}{ylabel}\PY{p}{(}\PY{l+s+sa}{r}\PY{l+s+s2}{\PYZdq{}}\PY{l+s+s2}{\PYZdl{} }\PY{l+s+s2}{\PYZbs{}}\PY{l+s+s2}{lambda(x) \PYZdl{}}\PY{l+s+s2}{\PYZdq{}}\PY{p}{)}
\PY{n}{plt}\PY{o}{.}\PY{n}{xlabel}\PY{p}{(}\PY{l+s+s2}{\PYZdq{}}\PY{l+s+s2}{x (pixels)}\PY{l+s+s2}{\PYZdq{}}\PY{p}{)}
\end{Verbatim}
\end{tcolorbox}

            \begin{tcolorbox}[breakable, size=fbox, boxrule=.5pt, pad at break*=1mm, opacityfill=0]
\prompt{Out}{outcolor}{242}{\boxspacing}
\begin{Verbatim}[commandchars=\\\{\}]
Text(0.5, 0, 'x (pixels)')
\end{Verbatim}
\end{tcolorbox}
        
    \begin{center}
    \adjustimage{max size={0.9\linewidth}{0.9\paperheight}}{output_98_1.png}
    \end{center}
    { \hspace*{\fill} \\}
    
    \textbf{\emph{Neón}}

    \begin{tcolorbox}[breakable, size=fbox, boxrule=1pt, pad at break*=1mm,colback=cellbackground, colframe=cellborder]
\prompt{In}{incolor}{243}{\boxspacing}
\begin{Verbatim}[commandchars=\\\{\}]
\PY{n}{plt}\PY{o}{.}\PY{n}{plot}\PY{p}{(}\PY{n}{improved\PYZus{}xval\PYZus{}guesses\PYZus{}ne}\PY{p}{,} \PY{n}{guessed\PYZus{}wavelengths\PYZus{}ne}\PY{p}{,} \PY{l+s+s1}{\PYZsq{}}\PY{l+s+s1}{o}\PY{l+s+s1}{\PYZsq{}}\PY{p}{)}
\PY{n}{plt}\PY{o}{.}\PY{n}{plot}\PY{p}{(}\PY{n}{xaxis\PYZus{}ne}\PY{p}{,} \PY{n}{wavelengths\PYZus{}ne}\PY{p}{,} \PY{l+s+s1}{\PYZsq{}}\PY{l+s+s1}{\PYZhy{}}\PY{l+s+s1}{\PYZsq{}}\PY{p}{)}
\PY{n}{plt}\PY{o}{.}\PY{n}{ylabel}\PY{p}{(}\PY{l+s+sa}{r}\PY{l+s+s2}{\PYZdq{}}\PY{l+s+s2}{\PYZdl{} }\PY{l+s+s2}{\PYZbs{}}\PY{l+s+s2}{lambda(x) \PYZdl{}}\PY{l+s+s2}{\PYZdq{}}\PY{p}{)}
\PY{n}{plt}\PY{o}{.}\PY{n}{xlabel}\PY{p}{(}\PY{l+s+s2}{\PYZdq{}}\PY{l+s+s2}{x (pixels)}\PY{l+s+s2}{\PYZdq{}}\PY{p}{)}
\end{Verbatim}
\end{tcolorbox}

            \begin{tcolorbox}[breakable, size=fbox, boxrule=.5pt, pad at break*=1mm, opacityfill=0]
\prompt{Out}{outcolor}{243}{\boxspacing}
\begin{Verbatim}[commandchars=\\\{\}]
Text(0.5, 0, 'x (pixels)')
\end{Verbatim}
\end{tcolorbox}
        
    \begin{center}
    \adjustimage{max size={0.9\linewidth}{0.9\paperheight}}{output_100_1.png}
    \end{center}
    { \hspace*{\fill} \\}
    
    \textbf{\emph{Kriptón}}

    \begin{tcolorbox}[breakable, size=fbox, boxrule=1pt, pad at break*=1mm,colback=cellbackground, colframe=cellborder]
\prompt{In}{incolor}{244}{\boxspacing}
\begin{Verbatim}[commandchars=\\\{\}]
\PY{n}{plt}\PY{o}{.}\PY{n}{plot}\PY{p}{(}\PY{n}{improved\PYZus{}xval\PYZus{}guesses\PYZus{}kr}\PY{p}{,} \PY{n}{guessed\PYZus{}wavelengths\PYZus{}kr}\PY{p}{,} \PY{l+s+s1}{\PYZsq{}}\PY{l+s+s1}{o}\PY{l+s+s1}{\PYZsq{}}\PY{p}{)}
\PY{n}{plt}\PY{o}{.}\PY{n}{plot}\PY{p}{(}\PY{n}{xaxis\PYZus{}kr}\PY{p}{,} \PY{n}{wavelengths\PYZus{}kr}\PY{p}{,} \PY{l+s+s1}{\PYZsq{}}\PY{l+s+s1}{\PYZhy{}}\PY{l+s+s1}{\PYZsq{}}\PY{p}{)}
\PY{n}{plt}\PY{o}{.}\PY{n}{ylabel}\PY{p}{(}\PY{l+s+sa}{r}\PY{l+s+s2}{\PYZdq{}}\PY{l+s+s2}{\PYZdl{} }\PY{l+s+s2}{\PYZbs{}}\PY{l+s+s2}{lambda(x) \PYZdl{}}\PY{l+s+s2}{\PYZdq{}}\PY{p}{)}
\PY{n}{plt}\PY{o}{.}\PY{n}{xlabel}\PY{p}{(}\PY{l+s+s2}{\PYZdq{}}\PY{l+s+s2}{x (pixels)}\PY{l+s+s2}{\PYZdq{}}\PY{p}{)}
\end{Verbatim}
\end{tcolorbox}

            \begin{tcolorbox}[breakable, size=fbox, boxrule=.5pt, pad at break*=1mm, opacityfill=0]
\prompt{Out}{outcolor}{244}{\boxspacing}
\begin{Verbatim}[commandchars=\\\{\}]
Text(0.5, 0, 'x (pixels)')
\end{Verbatim}
\end{tcolorbox}
        
    \begin{center}
    \adjustimage{max size={0.9\linewidth}{0.9\paperheight}}{output_102_1.png}
    \end{center}
    { \hspace*{\fill} \\}
    
    \textbf{\emph{Helio}}

    \begin{tcolorbox}[breakable, size=fbox, boxrule=1pt, pad at break*=1mm,colback=cellbackground, colframe=cellborder]
\prompt{In}{incolor}{245}{\boxspacing}
\begin{Verbatim}[commandchars=\\\{\}]
\PY{n}{plt}\PY{o}{.}\PY{n}{plot}\PY{p}{(}\PY{n}{improved\PYZus{}xval\PYZus{}guesses\PYZus{}he}\PY{p}{,} \PY{n}{guessed\PYZus{}wavelengths\PYZus{}he}\PY{p}{,} \PY{l+s+s1}{\PYZsq{}}\PY{l+s+s1}{o}\PY{l+s+s1}{\PYZsq{}}\PY{p}{)}
\PY{n}{plt}\PY{o}{.}\PY{n}{plot}\PY{p}{(}\PY{n}{xaxis\PYZus{}he}\PY{p}{,} \PY{n}{wavelengths\PYZus{}he}\PY{p}{,} \PY{l+s+s1}{\PYZsq{}}\PY{l+s+s1}{\PYZhy{}}\PY{l+s+s1}{\PYZsq{}}\PY{p}{)}
\PY{n}{plt}\PY{o}{.}\PY{n}{ylabel}\PY{p}{(}\PY{l+s+sa}{r}\PY{l+s+s2}{\PYZdq{}}\PY{l+s+s2}{\PYZdl{} }\PY{l+s+s2}{\PYZbs{}}\PY{l+s+s2}{lambda(x) \PYZdl{}}\PY{l+s+s2}{\PYZdq{}}\PY{p}{)}
\PY{n}{plt}\PY{o}{.}\PY{n}{xlabel}\PY{p}{(}\PY{l+s+s2}{\PYZdq{}}\PY{l+s+s2}{x (pixels)}\PY{l+s+s2}{\PYZdq{}}\PY{p}{)}
\end{Verbatim}
\end{tcolorbox}

            \begin{tcolorbox}[breakable, size=fbox, boxrule=.5pt, pad at break*=1mm, opacityfill=0]
\prompt{Out}{outcolor}{245}{\boxspacing}
\begin{Verbatim}[commandchars=\\\{\}]
Text(0.5, 0, 'x (pixels)')
\end{Verbatim}
\end{tcolorbox}
        
    \begin{center}
    \adjustimage{max size={0.9\linewidth}{0.9\paperheight}}{output_104_1.png}
    \end{center}
    { \hspace*{\fill} \\}
    
    ¡En efecto, un modelo lineal ajustó excelentemente!

    Podemos consultar las listas de líneas de neón, kriptón y helio para ver
qué tan bien lo hicimos y quizá intentar mejorar nuestro ajuste.

NIST, el Instituto Nacional de Estándares y Tecnología, mantiene
\href{https://physics.nist.gov/PhysRefData/ASD/lines_form.html}{listas
de líneas atómicas}. Sin embargo, ¡no tenemos una manera directa de
saber qué transiciones de estos átomos usar! Aun así, podemos aplicar
algunas heurísticas.

    \begin{tcolorbox}[breakable, size=fbox, boxrule=1pt, pad at break*=1mm,colback=cellbackground, colframe=cellborder]
\prompt{In}{incolor}{250}{\boxspacing}
\begin{Verbatim}[commandchars=\\\{\}]
\PY{c+c1}{\PYZsh{} we adopt the minimum/maximum wavelength from our linear fit}
\PY{n}{minwave\PYZus{}hg} \PY{o}{=} \PY{n}{wavelengths\PYZus{}hg}\PY{o}{.}\PY{n}{min}\PY{p}{(}\PY{p}{)}
\PY{n}{maxwave\PYZus{}hg} \PY{o}{=} \PY{n}{wavelengths\PYZus{}hg}\PY{o}{.}\PY{n}{max}\PY{p}{(}\PY{p}{)}

\PY{n}{minwave\PYZus{}ne} \PY{o}{=} \PY{n}{wavelengths\PYZus{}ne}\PY{o}{.}\PY{n}{min}\PY{p}{(}\PY{p}{)}
\PY{n}{maxwave\PYZus{}ne} \PY{o}{=} \PY{n}{wavelengths\PYZus{}ne}\PY{o}{.}\PY{n}{max}\PY{p}{(}\PY{p}{)}

\PY{n}{minwave\PYZus{}kr} \PY{o}{=} \PY{n}{wavelengths\PYZus{}kr}\PY{o}{.}\PY{n}{min}\PY{p}{(}\PY{p}{)}
\PY{n}{maxwave\PYZus{}kr} \PY{o}{=} \PY{n}{wavelengths\PYZus{}kr}\PY{o}{.}\PY{n}{max}\PY{p}{(}\PY{p}{)}

\PY{n}{minwave\PYZus{}he} \PY{o}{=} \PY{n}{wavelengths\PYZus{}he}\PY{o}{.}\PY{n}{min}\PY{p}{(}\PY{p}{)}
\PY{n}{maxwave\PYZus{}he} \PY{o}{=} \PY{n}{wavelengths\PYZus{}he}\PY{o}{.}\PY{n}{max}\PY{p}{(}\PY{p}{)}

\PY{c+c1}{\PYZsh{} then we search for atomic lines}
\PY{c+c1}{\PYZsh{} We are only interested in neutral lines, assuming the lamps are not hot enough to ionize the atoms}
\PY{n}{mercury\PYZus{}lines} \PY{o}{=} \PY{n}{Nist}\PY{o}{.}\PY{n}{query}\PY{p}{(}\PY{n}{minwav}\PY{o}{=}\PY{n}{minwave\PYZus{}hg}\PY{p}{,}
                           \PY{n}{maxwav}\PY{o}{=}\PY{n}{maxwave\PYZus{}hg}\PY{p}{,}
                           \PY{n}{linename}\PY{o}{=}\PY{l+s+s1}{\PYZsq{}}\PY{l+s+s1}{Hg I}\PY{l+s+s1}{\PYZsq{}}\PY{p}{)}
\PY{n}{krypton\PYZus{}lines} \PY{o}{=} \PY{n}{Nist}\PY{o}{.}\PY{n}{query}\PY{p}{(}\PY{n}{minwav}\PY{o}{=}\PY{n}{minwave\PYZus{}kr}\PY{p}{,}
                           \PY{n}{maxwav}\PY{o}{=}\PY{n}{maxwave\PYZus{}kr}\PY{p}{,}
                           \PY{n}{linename}\PY{o}{=}\PY{l+s+s1}{\PYZsq{}}\PY{l+s+s1}{Kr I}\PY{l+s+s1}{\PYZsq{}}\PY{p}{)}
\PY{n}{neon\PYZus{}lines} \PY{o}{=} \PY{n}{Nist}\PY{o}{.}\PY{n}{query}\PY{p}{(}\PY{n}{minwav}\PY{o}{=}\PY{n}{minwave\PYZus{}ne}\PY{p}{,}
                        \PY{n}{maxwav}\PY{o}{=}\PY{n}{maxwave\PYZus{}ne}\PY{p}{,}
                        \PY{n}{linename}\PY{o}{=}\PY{l+s+s1}{\PYZsq{}}\PY{l+s+s1}{Ne I}\PY{l+s+s1}{\PYZsq{}}\PY{p}{)}
\PY{n}{helium\PYZus{}lines} \PY{o}{=} \PY{n}{Nist}\PY{o}{.}\PY{n}{query}\PY{p}{(}\PY{n}{minwav}\PY{o}{=}\PY{n}{minwave\PYZus{}he}\PY{p}{,}
                         \PY{n}{maxwav}\PY{o}{=}\PY{n}{maxwave\PYZus{}he}\PY{p}{,}
                         \PY{n}{linename}\PY{o}{=}\PY{l+s+s1}{\PYZsq{}}\PY{l+s+s1}{He I}\PY{l+s+s1}{\PYZsq{}}\PY{p}{)}
\end{Verbatim}
\end{tcolorbox}

    Podemos comenzar mirando lo que NIST reporta para el espectro de
mercurio que ya ajustamos:

(nota que las longitudes de onda son longitudes de onda \emph{en aire},
no en vacío)

    \textbf{\emph{Mercurio}}

    \begin{tcolorbox}[breakable, size=fbox, boxrule=1pt, pad at break*=1mm,colback=cellbackground, colframe=cellborder]
\prompt{In}{incolor}{251}{\boxspacing}
\begin{Verbatim}[commandchars=\\\{\}]
\PY{n}{plt}\PY{o}{.}\PY{n}{figure}\PY{p}{(}\PY{n}{figsize}\PY{o}{=}\PY{p}{(}\PY{l+m+mi}{10}\PY{p}{,}\PY{l+m+mi}{4}\PY{p}{)}\PY{p}{)}
\PY{n}{plt}\PY{o}{.}\PY{n}{plot}\PY{p}{(}\PY{n}{wavelengths\PYZus{}hg}\PY{p}{,} \PY{n}{hg\PYZus{}spectrum}\PY{p}{,} \PY{n}{color}\PY{o}{=}\PY{l+s+s1}{\PYZsq{}}\PY{l+s+s1}{deepskyblue}\PY{l+s+s1}{\PYZsq{}}\PY{p}{)}
\PY{n}{plt}\PY{o}{.}\PY{n}{vlines}\PY{p}{(}\PY{n}{mercury\PYZus{}lines}\PY{p}{[}\PY{l+s+s1}{\PYZsq{}}\PY{l+s+s1}{Observed}\PY{l+s+s1}{\PYZsq{}}\PY{p}{]}\PY{p}{,} \PY{l+m+mi}{0}\PY{p}{,} \PY{n+nb}{max}\PY{p}{(}\PY{n}{hg\PYZus{}spectrum}\PY{p}{)}\PY{o}{*}\PY{l+m+mi}{1}\PY{p}{,} 
           \PY{n}{color}\PY{o}{=}\PY{l+s+s1}{\PYZsq{}}\PY{l+s+s1}{white}\PY{l+s+s1}{\PYZsq{}}\PY{p}{,} \PY{n}{alpha}\PY{o}{=}\PY{l+m+mf}{0.3}\PY{p}{,} \PY{n}{lw}\PY{o}{=}\PY{l+m+mf}{1.2}\PY{p}{)}
\PY{n}{plt}\PY{o}{.}\PY{n}{xlabel}\PY{p}{(}\PY{l+s+sa}{r}\PY{l+s+s2}{\PYZdq{}}\PY{l+s+s2}{Wavelength \PYZdl{}}\PY{l+s+s2}{\PYZbs{}}\PY{l+s+s2}{lambda\PYZdl{} (Å)}\PY{l+s+s2}{\PYZdq{}}\PY{p}{)}
\PY{n}{plt}\PY{o}{.}\PY{n}{ylabel}\PY{p}{(}\PY{l+s+s2}{\PYZdq{}}\PY{l+s+s2}{Intensity (a.u.)}\PY{l+s+s2}{\PYZdq{}}\PY{p}{)}
\PY{n}{plt}\PY{o}{.}\PY{n}{title}\PY{p}{(}\PY{l+s+s2}{\PYZdq{}}\PY{l+s+s2}{Mercury Lamp Spectrum and NIST Reference Lines}\PY{l+s+s2}{\PYZdq{}}\PY{p}{)}
\PY{n}{plt}\PY{o}{.}\PY{n}{show}\PY{p}{(}\PY{p}{)}
\end{Verbatim}
\end{tcolorbox}

    \begin{center}
    \adjustimage{max size={0.9\linewidth}{0.9\paperheight}}{output_110_0.png}
    \end{center}
    { \hspace*{\fill} \\}
    
    \textbf{\emph{Neón}}

    \begin{tcolorbox}[breakable, size=fbox, boxrule=1pt, pad at break*=1mm,colback=cellbackground, colframe=cellborder]
\prompt{In}{incolor}{252}{\boxspacing}
\begin{Verbatim}[commandchars=\\\{\}]
\PY{n}{plt}\PY{o}{.}\PY{n}{figure}\PY{p}{(}\PY{n}{figsize}\PY{o}{=}\PY{p}{(}\PY{l+m+mi}{10}\PY{p}{,}\PY{l+m+mi}{4}\PY{p}{)}\PY{p}{)}
\PY{n}{plt}\PY{o}{.}\PY{n}{plot}\PY{p}{(}\PY{n}{wavelengths\PYZus{}ne}\PY{p}{,} \PY{n}{ne\PYZus{}spectrum}\PY{p}{,} \PY{n}{color}\PY{o}{=}\PY{l+s+s1}{\PYZsq{}}\PY{l+s+s1}{deepskyblue}\PY{l+s+s1}{\PYZsq{}}\PY{p}{)}
\PY{n}{plt}\PY{o}{.}\PY{n}{vlines}\PY{p}{(}\PY{n}{neon\PYZus{}lines}\PY{p}{[}\PY{l+s+s1}{\PYZsq{}}\PY{l+s+s1}{Observed}\PY{l+s+s1}{\PYZsq{}}\PY{p}{]}\PY{p}{,} \PY{l+m+mi}{0}\PY{p}{,} \PY{n+nb}{max}\PY{p}{(}\PY{n}{ne\PYZus{}spectrum}\PY{p}{)}\PY{o}{*}\PY{l+m+mi}{1}\PY{p}{,} 
           \PY{n}{color}\PY{o}{=}\PY{l+s+s1}{\PYZsq{}}\PY{l+s+s1}{white}\PY{l+s+s1}{\PYZsq{}}\PY{p}{,} \PY{n}{alpha}\PY{o}{=}\PY{l+m+mf}{0.3}\PY{p}{,} \PY{n}{lw}\PY{o}{=}\PY{l+m+mf}{1.2}\PY{p}{)}
\PY{n}{plt}\PY{o}{.}\PY{n}{xlabel}\PY{p}{(}\PY{l+s+sa}{r}\PY{l+s+s2}{\PYZdq{}}\PY{l+s+s2}{Wavelength \PYZdl{}}\PY{l+s+s2}{\PYZbs{}}\PY{l+s+s2}{lambda\PYZdl{} (Å)}\PY{l+s+s2}{\PYZdq{}}\PY{p}{)}
\PY{n}{plt}\PY{o}{.}\PY{n}{ylabel}\PY{p}{(}\PY{l+s+s2}{\PYZdq{}}\PY{l+s+s2}{Intensity (a.u.)}\PY{l+s+s2}{\PYZdq{}}\PY{p}{)}
\PY{n}{plt}\PY{o}{.}\PY{n}{title}\PY{p}{(}\PY{l+s+s2}{\PYZdq{}}\PY{l+s+s2}{Neon Lamp Spectrum and NIST Reference Lines}\PY{l+s+s2}{\PYZdq{}}\PY{p}{)}
\PY{n}{plt}\PY{o}{.}\PY{n}{show}\PY{p}{(}\PY{p}{)}
\end{Verbatim}
\end{tcolorbox}

    \begin{center}
    \adjustimage{max size={0.9\linewidth}{0.9\paperheight}}{output_112_0.png}
    \end{center}
    { \hspace*{\fill} \\}
    
    \textbf{\emph{Kriptón}}

    \begin{tcolorbox}[breakable, size=fbox, boxrule=1pt, pad at break*=1mm,colback=cellbackground, colframe=cellborder]
\prompt{In}{incolor}{253}{\boxspacing}
\begin{Verbatim}[commandchars=\\\{\}]
\PY{n}{plt}\PY{o}{.}\PY{n}{figure}\PY{p}{(}\PY{n}{figsize}\PY{o}{=}\PY{p}{(}\PY{l+m+mi}{10}\PY{p}{,}\PY{l+m+mi}{4}\PY{p}{)}\PY{p}{)}
\PY{n}{plt}\PY{o}{.}\PY{n}{plot}\PY{p}{(}\PY{n}{wavelengths\PYZus{}kr}\PY{p}{,} \PY{n}{kr\PYZus{}spectrum}\PY{p}{,} \PY{n}{color}\PY{o}{=}\PY{l+s+s1}{\PYZsq{}}\PY{l+s+s1}{deepskyblue}\PY{l+s+s1}{\PYZsq{}}\PY{p}{)}
\PY{n}{plt}\PY{o}{.}\PY{n}{vlines}\PY{p}{(}\PY{n}{krypton\PYZus{}lines}\PY{p}{[}\PY{l+s+s1}{\PYZsq{}}\PY{l+s+s1}{Observed}\PY{l+s+s1}{\PYZsq{}}\PY{p}{]}\PY{p}{,} \PY{l+m+mi}{0}\PY{p}{,} \PY{n+nb}{max}\PY{p}{(}\PY{n}{kr\PYZus{}spectrum}\PY{p}{)}\PY{o}{*}\PY{l+m+mi}{1}\PY{p}{,} 
           \PY{n}{color}\PY{o}{=}\PY{l+s+s1}{\PYZsq{}}\PY{l+s+s1}{white}\PY{l+s+s1}{\PYZsq{}}\PY{p}{,} \PY{n}{alpha}\PY{o}{=}\PY{l+m+mf}{0.3}\PY{p}{,} \PY{n}{lw}\PY{o}{=}\PY{l+m+mf}{1.2}\PY{p}{)}
\PY{n}{plt}\PY{o}{.}\PY{n}{xlabel}\PY{p}{(}\PY{l+s+sa}{r}\PY{l+s+s2}{\PYZdq{}}\PY{l+s+s2}{Wavelength \PYZdl{}}\PY{l+s+s2}{\PYZbs{}}\PY{l+s+s2}{lambda\PYZdl{} (Å)}\PY{l+s+s2}{\PYZdq{}}\PY{p}{)}
\PY{n}{plt}\PY{o}{.}\PY{n}{ylabel}\PY{p}{(}\PY{l+s+s2}{\PYZdq{}}\PY{l+s+s2}{Intensity (a.u.)}\PY{l+s+s2}{\PYZdq{}}\PY{p}{)}
\PY{n}{plt}\PY{o}{.}\PY{n}{title}\PY{p}{(}\PY{l+s+s2}{\PYZdq{}}\PY{l+s+s2}{Krypton Lamp Spectrum and NIST Reference Lines}\PY{l+s+s2}{\PYZdq{}}\PY{p}{)}
\PY{n}{plt}\PY{o}{.}\PY{n}{show}\PY{p}{(}\PY{p}{)}
\end{Verbatim}
\end{tcolorbox}

    \begin{center}
    \adjustimage{max size={0.9\linewidth}{0.9\paperheight}}{output_114_0.png}
    \end{center}
    { \hspace*{\fill} \\}
    
    \textbf{\emph{Helio}}

    \begin{tcolorbox}[breakable, size=fbox, boxrule=1pt, pad at break*=1mm,colback=cellbackground, colframe=cellborder]
\prompt{In}{incolor}{254}{\boxspacing}
\begin{Verbatim}[commandchars=\\\{\}]
\PY{n}{plt}\PY{o}{.}\PY{n}{figure}\PY{p}{(}\PY{n}{figsize}\PY{o}{=}\PY{p}{(}\PY{l+m+mi}{10}\PY{p}{,}\PY{l+m+mi}{4}\PY{p}{)}\PY{p}{)}
\PY{n}{plt}\PY{o}{.}\PY{n}{plot}\PY{p}{(}\PY{n}{wavelengths\PYZus{}he}\PY{p}{,} \PY{n}{he\PYZus{}spectrum}\PY{p}{,} \PY{n}{color}\PY{o}{=}\PY{l+s+s1}{\PYZsq{}}\PY{l+s+s1}{deepskyblue}\PY{l+s+s1}{\PYZsq{}}\PY{p}{)}
\PY{n}{plt}\PY{o}{.}\PY{n}{vlines}\PY{p}{(}\PY{n}{helium\PYZus{}lines}\PY{p}{[}\PY{l+s+s1}{\PYZsq{}}\PY{l+s+s1}{Observed}\PY{l+s+s1}{\PYZsq{}}\PY{p}{]}\PY{p}{,} \PY{l+m+mi}{0}\PY{p}{,} \PY{n+nb}{max}\PY{p}{(}\PY{n}{he\PYZus{}spectrum}\PY{p}{)}\PY{o}{*}\PY{l+m+mi}{1}\PY{p}{,} 
           \PY{n}{color}\PY{o}{=}\PY{l+s+s1}{\PYZsq{}}\PY{l+s+s1}{white}\PY{l+s+s1}{\PYZsq{}}\PY{p}{,} \PY{n}{alpha}\PY{o}{=}\PY{l+m+mf}{0.3}\PY{p}{,} \PY{n}{lw}\PY{o}{=}\PY{l+m+mf}{1.2}\PY{p}{)}
\PY{n}{plt}\PY{o}{.}\PY{n}{xlabel}\PY{p}{(}\PY{l+s+sa}{r}\PY{l+s+s2}{\PYZdq{}}\PY{l+s+s2}{Wavelength \PYZdl{}}\PY{l+s+s2}{\PYZbs{}}\PY{l+s+s2}{lambda\PYZdl{} (Å)}\PY{l+s+s2}{\PYZdq{}}\PY{p}{)}
\PY{n}{plt}\PY{o}{.}\PY{n}{ylabel}\PY{p}{(}\PY{l+s+s2}{\PYZdq{}}\PY{l+s+s2}{Intensity (a.u.)}\PY{l+s+s2}{\PYZdq{}}\PY{p}{)}
\PY{n}{plt}\PY{o}{.}\PY{n}{title}\PY{p}{(}\PY{l+s+s2}{\PYZdq{}}\PY{l+s+s2}{Helium Lamp Spectrum and NIST Reference Lines}\PY{l+s+s2}{\PYZdq{}}\PY{p}{)}
\PY{n}{plt}\PY{o}{.}\PY{n}{show}\PY{p}{(}\PY{p}{)}
\end{Verbatim}
\end{tcolorbox}

    \begin{center}
    \adjustimage{max size={0.9\linewidth}{0.9\paperheight}}{output_116_0.png}
    \end{center}
    { \hspace*{\fill} \\}
    
    ¡Has terminado! Ahora aplica esta calibración a los datos y comienza a
medir cosas en el Tutorial 3.


    % Add a bibliography block to the postdoc
    
    
    
\end{document}
